\chapter[Latent Lie-Algebra Links]{Latent Lie-Algebra Links in the Gaussian Realization}
\label{ch:gauss-latent-links}

What do the latent exponential-chart links of \Cref{ch:links} become when every belief law is
Gaussian, and what does Gaussian matter do to them when it is integrated out? In this realization
the law-level transported divergence is an explicit function of the generator whose log-volume
part is linear in it (\Cref{prop:glink-law-level-energy}). The expected energy of a linear
comparison record gives the delta-rule current, and without a link prior its minimizing transport
matrix has rank at most one and, for \(K\geq2\), lies outside the exponential chart, while stationary
points in the chart can still sit at critical points of the exponential map
(\Cref{prop:glink-sample-record}); when its trace is negative, no chart link with semisimple negative
eigenvalues attains the infimum over the chart, so for \(K=2\) and \(K=3\) the infimum is not attained
(\Cref{prop:glink-chart-infimum-semisimple}). Under a product Gaussian recognition law of the
generators the expected holonomy of a closed walk that traverses each reversal pair at most once is
an ordered product of per-copy exponential moments
(\Cref{prop:glink-expected-loop}). Jointly normalized undirected Gaussian matter induces on the
generators the action \(\mathscr S^{\mathrm{ind}}=-\mathsf A(\mathsf h,\Lambda(\mathring\omega))\),
one half the log-determinant of a connection-Laplacian precision at zero information vector,
bounded below by the self terms (\Cref{prop:glink-induced-action}); its hopping expansion is a
convergent sum of weighted Wilson loops whose leading link-dependent term on a triangle of agents,
for isotropic data and isometric links, is a multiple of the loop defect of the triangle
(\Cref{prop:glink-hopping}). The link-dependent normalizers of
\Cref{prop:cg-normalized-schur-action,prop:obs-normalizer-link-dependence,prop:obs-holonomy-determinant-factorization}
are values of this induced action at fixed links, and an unanchored cycle, whose law at fixed links
exists off a null set, admits no joint law with latent links (\Cref{cor:glink-unanchored-cycle}),
while for \(K=1\) unanchored matter on a connected multigraph of cycle rank at least two has one
(\Cref{prop:glink-unanchored-scalar}). An explicit spectrum of the linearized link operator on a
ring closes the chapter (\Cref{prop:glink-linear-spectrum}). These are the Gaussian values of
\Cref{prop:link-law-level-energy,cor:link-delta-rule,cor:link-prior-necessary,prop:link-induced-action,prop:link-linearized-equation}
of the main manuscript, and the induced action is the quantity against which a loop action of the
generator prior of \Cref{def:link-generator-prior} can be compared.

Every statement concerns the Gaussian realization with the specialization
\(\Gamma_b=\GL^+(K)\) of \Cref{hyp:gen-lg-frame-group}, one channel, and held model latents. The
chapter does not claim that inference of the generators reaches the stationary points it computes,
does not select anchors, generator covariances, loop coefficients, or increment precisions, does
not identify the induced action with a loop action of the generator prior beyond the leading
hopping order, gives no expansion of the induced action over cycle-rooted spanning forests in rank
three or more, where a combinatorial interpretation is an open question of
\citet[Sec.~14]{Kenyon2011VectorBundleLaplacian}, and evaluates expected loop defects as ordered
products only for walks without a repeated reversal pair; the ring of \Cref{sec:glink-spectrum} illustrates
the linearized operator and sets no direction for the theory. \status{NOT-CLAIMED} Joint
normalization of unanchored matter on multigraphs of cycle rank at least two at a stated
\(K\geq2\) is \Cref{open:glink-unanchored-normalization}, and attainment of the chart infimum of an
unregularized record energy for \(K\geq4\) at a chart link with a non-semisimple negative eigenvalue
is \Cref{open:glink-chart-infimum}.

\section{Generator coordinates and reference laws}
\label{sec:glink-setting}

Throughout the chapter \(G=\GL^+(K)\) with \(\rho_b\) its defining representation, which realizes
\(\Gamma_b=\GL^+(K)\) (\Cref{hyp:gen-lg-frame-group}); the Lie algebra is
\(\mathfrak g=\mathfrak{gl}(K,\R)\), \(\mathrm d\rho_b\) is the identity, and the exponential is
the matrix exponential. Only the belief channel is considered, and its index is dropped. Each
oriented edge copy \(e:j\to i\) carries a generator \(\mathring\omega_e\in\mathfrak g\) with twist
\(\mathring\Theta_e=\exp(\mathring\omega_e)\) and \(\mathring\omega_{\bar e}=-\mathring\omega_e\)
(\Cref{def:link-exponential-chart}); \(E_+\) contains one orientation of each reversal pair and
\(\mathscr T=\mathfrak g^{E_+}\). A closed walk \(p=(e_0,\ldots,e_{r-1})\) with
\(e_a:i_{a+1}\to i_a\) and \(i_r=i_0\) has the reference-coordinate holonomy
\(\mathring\Theta_p=\mathring\Theta_{e_0}\cdots\mathring\Theta_{e_{r-1}}\), and
\(\mathrm{Hol}(p)=U_{i_0}\mathring\Theta_pU_{i_0}^{-1}\) by \eqref{eq:geo-edge-twist-holonomy}.

Laws, self terms, weights, observations, and information vectors are written in reference
coordinates: the belief law of agent \(i\) stands for \(\rho_b(U_i)^{-1}_\#q_i\). In these
coordinates the link of \(e:j\to i\) acts as \(\exp(\mathring\omega_e)\), because the passive
re-expression \(g_i=U_i^{-1}\) makes every frame the reference frame, leaves every twist unchanged
(\Cref{prop:geo-edge-twist-invariance}), and pushes both arguments of every transported divergence
forward by one bijection, which preserves relative entropy. A Gaussian energy in local coordinates
becomes one in reference coordinates by the block congruence of
\Cref{rem:gauss-membership-frame-dependence}, which changes a log-determinant by the
link-independent constant \(2\sum_i\log\lvert\det U_i\rvert\), the Jacobian bookkeeping of
\Cref{prop:cg-nonflat-series-gauge}.

Matrices are paired by the trace pairing \(\langle\mathbf S,\mathbf T\rangle=\Tr(\mathbf S^\top\mathbf T)\),
and gradients are taken for this pairing; the gradient for the inner product
\(\langle\cdot,\cdot\rangle_{\mathsf G}\) of \Cref{ch:links} is \(\mathsf G^{-1}\nabla\mathsf G\), with
\(\mathsf G=\mathsf G^{\mathrm{fib}}\). Kronecker products have the coordinates
\((\mathbf S\otimes\mathbf T)_{(a,b),(c,d)}=\mathbf S_{ac}\mathbf T_{bd}\). The Fr\'echet derivative of the
exponential is \(\mathbf L(\mathbf B,\mathbf T)=\int_0^1e^{t\mathbf B}\mathbf Te^{(1-t)\mathbf B}dt\)
(\Cref{lem:link-exp-derivative}), and cyclicity of the trace under the integral makes
\(\mathbf L(\mathbf B^\top,\cdot)\) its adjoint for the trace pairing. The second derivative
\(\mathbf L^{(2)}(\mathbf B;\mathbf T)=\frac{d^2}{dt^2}e^{\mathbf B+t\mathbf T}\big|_{t=0}\) is used once. The loop
defect of a closed walk is
\(\mathfrak Q_p=\tfrac12\langle\mathring\Theta_p-I,\mathring\Theta_p-I\rangle_{\mathsf G}\)
(\Cref{def:link-loop-defect}). Bold capitals \(\mathbf B,\mathbf S,\mathbf T\) denote generic matrices
inside one statement or proof, and upright Roman letters denote statement-local real constants.

\section{Law-level and record energies of one copy}
\label{sec:glink-energies}

\propositionheading{Law-level transported divergence in generator coordinates}{prop:glink-law-level-energy}
Let \(q_i=\mathcal N(\mu_i,\Sigma_i)\) and \(q_j=\mathcal N(\mu_j,\Sigma_j)\) be nondegenerate
Gaussian laws in reference coordinates, \(e:j\to i\), and \(\mathring\omega=\mathring\omega_e\in\mathfrak g\).
(i) The transported divergence is
\begin{equation}
\begin{aligned}
\KL\bigl(q_i\Vert\exp(\mathring\omega)_\#q_j\bigr)
={}&\tfrac12\Tr\bigl(\Sigma_j^{-1}e^{-\mathring\omega}\Sigma_ie^{-\mathring\omega^\top}\bigr)
+\tfrac12\bigl(e^{-\mathring\omega}\mu_i-\mu_j\bigr)^\top\Sigma_j^{-1}\bigl(e^{-\mathring\omega}\mu_i-\mu_j\bigr)\\
&-\tfrac K2+\tfrac12\log\frac{\det\Sigma_j}{\det\Sigma_i}+\Tr\mathring\omega ,
\end{aligned}
\label{eq:glink-law-level-energy}
\end{equation}
and it equals \(\KL(q_i\Vert(\Theta_e)_\#q_j)\) computed in local coordinates for every choice of
frames. The log-volume part \(\Tr\mathring\omega\) is linear in the generator, and every other
dependence on \(\mathring\omega\) passes through \(e^{-\mathring\omega}\).
(ii) The divergence is a real-analytic function of \(\mathring\omega\) with differential and gradient
\begin{equation}
\begin{gathered}
d\KL[\mathbf T]=\Tr\mathbf T-\bigl\langle\mathbf A_e,\mathbf L(-\mathring\omega,\mathbf T)\bigr\rangle,
\qquad
\nabla_{\mathring\omega}\KL=I-\mathbf L(-\mathring\omega^\top,\mathbf A_e),\\
\mathbf A_e=\Sigma_j^{-1}\bigl(e^{-\mathring\omega}(\Sigma_i+\mu_i\mu_i^\top)-\mu_j\mu_i^\top\bigr).
\end{gathered}
\label{eq:glink-law-level-gradient}
\end{equation}
(iii) At \(\mathring\omega=0\) it is the Gaussian relative entropy of \(q_i\) to \(q_j\). If the
transported source has the receiver's covariance, \(\Sigma_i=e^{\mathring\omega}\Sigma_je^{\mathring\omega^\top}\),
it is the twisted pair energy \eqref{eq:cg-twisted-pair} with \(W_e=\Sigma_i^{-1}\),
\(\tfrac12(\mu_i-e^{\mathring\omega}\mu_j)^\top\Sigma_i^{-1}(\mu_i-e^{\mathring\omega}\mu_j)\), which
at \(\mathring\omega=0\) and \(\Sigma_i=\Sigma_j\) is the edge relative entropy of
\Cref{prop:ig-fixed-covariance-kl}. If \(q_i=q_j=\mathcal N(\mu,\Sigma)\) and
\(\mathring\omega=t\mathring\upsilon\), then as \(t\to0\) the divergence is
\(\tfrac{t^2}2\mathfrak k_q(\mathring\upsilon,\mathring\upsilon)+O(t^3)\) with the Gaussian orbit
Fisher form
\begin{equation}
\mathfrak k_q(\mathring\upsilon,\mathring\upsilon)
=(\mathring\upsilon\mu)^\top\Sigma^{-1}\mathring\upsilon\mu
+\Tr(\mathring\upsilon^2)+\Tr\bigl(\Sigma^{-1}\mathring\upsilon\Sigma\mathring\upsilon^\top\bigr)
\label{eq:glink-orbit-fisher}
\end{equation}
of \Cref{cor:pb-orbit-fisher-form}. \status{ESTABLISHED}

\noindent\emph{Proof.}
(i) By \Cref{prop:geo-moment-pushforward},
\(\exp(\mathring\omega)_\#q_j=\mathcal N(e^{\mathring\omega}\mu_j,e^{\mathring\omega}\Sigma_je^{\mathring\omega^\top})\).
For nondegenerate Gaussian laws on \(\R^K\), integrating the difference of the log densities gives
\(\KL(\mathcal N(\mu_i,\Sigma_i)\Vert\mathcal N(\mu',\Sigma'))=\tfrac12[\Tr(\Sigma'^{-1}\Sigma_i)+(\mu_i-\mu')^\top\Sigma'^{-1}(\mu_i-\mu')-K+\log\det\Sigma'-\log\det\Sigma_i]\).
With \(\mu'=e^{\mathring\omega}\mu_j\) and \(\Sigma'=e^{\mathring\omega}\Sigma_je^{\mathring\omega^\top}\), the inverse
is \(e^{-\mathring\omega^\top}\Sigma_j^{-1}e^{-\mathring\omega}\), cyclicity of the trace moves the
outer factors onto \(\Sigma_i\) and onto the mean difference, and
\(\log\det\Sigma'=\log\det\Sigma_j+2\log\det e^{\mathring\omega}\). Jacobi's formula
\(\frac{d}{dt}\det e^{t\mathring\omega}=\det e^{t\mathring\omega}\Tr(e^{-t\mathring\omega}\mathring\omega e^{t\mathring\omega})=\Tr(\mathring\omega)\det e^{t\mathring\omega}\)
integrates to \(\det e^{\mathring\omega}=e^{\Tr\mathring\omega}\), which gives
\eqref{eq:glink-law-level-energy}. In local coordinates the two laws are
\(\rho_b(U_i)_\#q_i\) and
\((U_ie^{\mathring\omega}U_j^{-1})_\#\rho_b(U_j)_\#q_j=\rho_b(U_i)_\#\exp(\mathring\omega)_\#q_j\),
and relative entropy is invariant under the common bijection \(\rho_b(U_i)\).
(ii) The right side of \eqref{eq:glink-law-level-energy} is a polynomial in the entries of
\(\mathring\Theta_e^{-1}=e^{-\mathring\omega}\) plus \(\Tr\mathring\omega\), and the exponential is entire. As
a function of \(\mathring\Theta_e^{-1}\), the first two terms have the differential
\(\langle\mathbf A_e,d\mathring\Theta_e^{-1}\rangle\): the first contributes
\(\langle\Sigma_j^{-1}\mathring\Theta_e^{-1}\Sigma_i,d\mathring\Theta_e^{-1}\rangle\) because \(\Sigma_i\) and
\(\Sigma_j\) are symmetric, and the second
\(\langle\Sigma_j^{-1}(\mathring\Theta_e^{-1}\mu_i-\mu_j)\mu_i^\top,d\mathring\Theta_e^{-1}\rangle\). By
\Cref{lem:link-exp-derivative}, the derivative of \(e^{-\mathring\omega}\) in the direction \(\mathbf T\) is
\(\mathbf L(-\mathring\omega,-\mathbf T)=-\mathbf L(-\mathring\omega,\mathbf T)\), and
\(d\Tr\mathring\omega[\mathbf T]=\Tr\mathbf T\), which gives the differential; the adjoint of
\(\mathbf L(-\mathring\omega,\cdot)\) is \(\mathbf L(-\mathring\omega^\top,\cdot)\), which gives the
gradient. (iii) The first two cases are substitution. For the third, insert
\(e^{-t\mathring\upsilon}=I-t\mathring\upsilon+\tfrac{t^2}2\mathring\upsilon^2+O(t^3)\) into
\eqref{eq:glink-law-level-energy} with \(\mu_i=\mu_j=\mu\) and \(\Sigma_i=\Sigma_j=\Sigma\). The trace term is
\(\tfrac K2-t\Tr\mathring\upsilon+\tfrac{t^2}2[\Tr(\mathring\upsilon^2)+\Tr(\Sigma^{-1}\mathring\upsilon\Sigma\mathring\upsilon^\top)]+O(t^3)\),
because the trace is invariant under cyclic permutation and transposition, so that the first-order
terms \(\Tr(\Sigma^{-1}\mathring\upsilon\Sigma)\) and \(\Tr\mathring\upsilon^\top\) both equal
\(\Tr\mathring\upsilon\), and the terms \(\Tr(\Sigma^{-1}\mathring\upsilon^2\Sigma)\) and
\(\Tr(\mathring\upsilon^\top)^2\) both equal \(\Tr(\mathring\upsilon^2)\); the mean term is
\(\tfrac{t^2}2(\mathring\upsilon\mu)^\top\Sigma^{-1}\mathring\upsilon\mu+O(t^3)\), since
\(e^{-t\mathring\upsilon}\mu-\mu=-t\mathring\upsilon\mu+O(t^2)\); the log-determinant term vanishes, and the
linear terms \(-t\Tr\mathring\upsilon\) and \(t\Tr\mathring\upsilon\) cancel. The sum is
\(\tfrac{t^2}2\) times \eqref{eq:glink-orbit-fisher} plus \(O(t^3)\). The form is the orbit Fisher form of
\Cref{cor:pb-orbit-fisher-form}: the moment curve
\(t\mapsto(e^{t\mathring\upsilon}\mu,e^{t\mathring\upsilon}\Sigma e^{t\mathring\upsilon^\top})\) has
velocity \((\mathring\upsilon\mu,\mathring\upsilon\Sigma+\Sigma\mathring\upsilon^\top)\) at \(t=0\),
the fundamental vector field of \(\mathring\upsilon\) at \(q\), and the Fisher metric
\eqref{eq:ig-moment-metric} evaluates it to
\((\mathring\upsilon\mu)^\top\Sigma^{-1}\mathring\upsilon\mu+\tfrac12\Tr[(\Sigma^{-1}\mathring\upsilon\Sigma+\mathring\upsilon^\top)^2]\);
expanding the square, the two pure terms each have trace \(\Tr(\mathring\upsilon^2)\) and the two
cross terms each \(\Tr(\Sigma^{-1}\mathring\upsilon\Sigma\mathring\upsilon^\top)\), which is
\eqref{eq:glink-orbit-fisher}.
\hfill$\square$

The divergence charges the generator in two ways: through \(e^{-\mathring\omega}\), which pulls the
receiver's law back into the source frame, and through the volume term \(\Tr\mathring\omega\), which
penalizes expansion of the source and rewards contraction. For trace-free generators, the
isometries of \Cref{sec:glink-induced} among them, the volume term vanishes.

The comparison record of \Cref{prop:obs-selective-gaussian-record}(iii) with \(S^Q_e=I\) and
\(S^K_e=\mathbf M\) is the kernel \(o_e\given z_i,z_j\sim\mathcal N(z_i-\mathbf Mz_j,W_e^{-1})\),
\(W_e\succ0\), a Gaussian kernel for every transport matrix \(\mathbf M\in\R^{K\times K}\), singular ones
included; its energy at an
observed \(o_e\) is
\begin{equation}
\mathsf E_e(z_i,z_j;\mathbf M)
=\tfrac12(o_e-z_i+\mathbf Mz_j)^\top W_e(o_e-z_i+\mathbf Mz_j)+\tfrac12\log\det(2\pi W_e^{-1}),
\label{eq:glink-record-energy}
\end{equation}
and the latent link sets \(\mathbf M=e^{\mathring\omega_e}\). At \(o_e=0\) its quadratic part is the record
energy of \Cref{cor:link-delta-rule}(a), and \(\mathbf M\) plays the role of the matrix of
\Cref{cor:link-prior-necessary}(a).

\propositionheading{Linear comparison record, delta-rule current, and the unregularized transport}{prop:glink-sample-record}
Let \(q_i=\mathcal N(\mu_i,\Sigma_i)\) and \(q_j=\mathcal N(\mu_j,\Sigma_j)\) be nondegenerate and
let matter be recognized by the product \(q_i\otimes q_j\). (i) The expected energy
\eqref{eq:glink-record-energy} of the comparison record with weight \(W_e\succ0\) at an observation
\(o_e\in\R^K\) is
\begin{equation}
\begin{aligned}
\E_{q_i\otimes q_j}\mathsf E_e(\cdot;\mathbf M)
={}&\tfrac12(\mu_i-o_e-\mathbf M\mu_j)^\top W_e(\mu_i-o_e-\mathbf M\mu_j)+\tfrac12\Tr(W_e\Sigma_i)\\
&+\tfrac12\Tr\bigl(W_e\mathbf M\Sigma_j\mathbf M^\top\bigr)+\tfrac12\log\det(2\pi W_e^{-1}),
\end{aligned}
\label{eq:glink-expected-record}
\end{equation}
and if \(\mathbf M=e^{\mathring\omega_e}\) with \(\mathring\omega_e\) random and independent of the
matter, its expectation over \(\mathring\omega_e\) is the affine function
\(\tfrac12(\mu_i-o_e)^\top W_e(\mu_i-o_e)-(\mu_i-o_e)^\top W_e(\E e^{\mathring\omega_e})\mu_j
+\tfrac12\sum_{a,b,c,d}(W_e)_{ab}(\Sigma_j+\mu_j\mu_j^\top)_{cd}\bigl(\E[e^{\mathring\omega_e}\otimes e^{\mathring\omega_e}]\bigr)_{(b,a),(c,d)}
+\tfrac12\Tr(W_e\Sigma_i)+\tfrac12\log\det(2\pi W_e^{-1})\)
of the first two exponential moments.
(ii) The derivative of \eqref{eq:glink-expected-record} in \(\mathbf M\) is \(-\mathbf J_e\), with the
current
\begin{equation}
\mathbf J_e=W_e\bigl[(\mu_i-o_e-\mathbf M\mu_j)\mu_j^\top-\mathbf M\Sigma_j\bigr],
\label{eq:glink-current}
\end{equation}
a precision-weighted prediction error times the source mean, minus the source covariance
transported by \(\mathbf M\), and its gradient in the generator at \(\mathbf M=e^{\mathring\omega_e}\)
is \(-\mathbf L(\mathring\omega_e^\top,\mathbf J_e)\); this is the Gaussian form of the current of
\Cref{cor:link-delta-rule} with an observed offset.
(iii) Without a link prior, \eqref{eq:glink-expected-record} is a strictly convex quadratic
function of \(\mathbf M\in\R^{K\times K}\) with the unique minimizer
\begin{equation}
\mathbf M^\star=(\mu_i-o_e)\mu_j^\top(\Sigma_j+\mu_j\mu_j^\top)^{-1},
\qquad
\Tr\mathbf M^\star=\frac{\mu_j^\top\Sigma_j^{-1}(\mu_i-o_e)}{1+\mu_j^\top\Sigma_j^{-1}\mu_j},
\label{eq:glink-rank-one}
\end{equation}
of rank at most one. For \(K\geq2\) it is singular and equals no \(e^{\mathring\omega}\);
consequently the gradient in \(\mathring\omega_e\) vanishes only where \(\mathbf L(\mathring\omega_e,\cdot)\)
is singular, that is, where two eigenvalues of \(\mathring\omega_e\) differ by a nonzero integer
multiple of \(2\pi\mathrm i\). For \(K\geq2\), if \(\Tr\mathbf M^\star\geq0\), the infimum over
\(\mathring\omega_e\in\mathfrak g\) equals the value at \(\mathbf M^\star\) and is not attained; if
\(\Tr\mathbf M^\star<0\), no sequence of exponentials converges to \(\mathbf M^\star\), and the
infimum over \(\mathfrak g\) exceeds the value at \(\mathbf M^\star\).
(iv) A stationary point at a critical point of the exponential can be a saddle. For \(K=2\),
\(W_e=I_2\), \(\Sigma_j=\varsigma^2I_2\), \(\mu_j=(1,0)^\top\), and
\(\mu_i-o_e=-(1+2\varsigma^2)e^{\mathrm a}(1,0)^\top\) with \(\varsigma>0\) and \(\mathrm a\in\R\), the
generator \(\mathring\omega_0=\mathrm aI_2+\pi\mathfrak J\), with
\(\mathfrak J=\begin{psmallmatrix}0&-1\\1&0\end{psmallmatrix}\), is a stationary point of
\(\mathring\omega\mapsto\E\mathsf E_e(\cdot;e^{\mathring\omega})\) with
\(e^{\mathring\omega_0}=-e^{\mathrm a}I_2\), and it is not a local minimizer.
\status{ESTABLISHED}

\noindent\emph{Proof.}
(i) Under \(q_i\otimes q_j\) the residual \(o_e-z_i+\mathbf Mz_j\) is Gaussian with mean
\(-(\mu_i-o_e-\mathbf M\mu_j)\) and covariance \(\Sigma_i+\mathbf M\Sigma_j\mathbf M^\top\), and the
expectation of a quadratic form is the form at the mean plus the trace against the covariance.
Expanding the first and third terms as
\(\tfrac12(\mu_i-o_e)^\top W_e(\mu_i-o_e)-(\mu_i-o_e)^\top W_e\mathbf M\mu_j+\tfrac12\Tr(W_e\mathbf M\mathbf S\mathbf M^\top)\)
with \(\mathbf S=\Sigma_j+\mu_j\mu_j^\top\), and writing
\(\Tr(W_e\mathbf M\mathbf S\mathbf M^\top)=\sum_{a,b,c,d}(W_e)_{ab}\mathbf M_{bc}\mathbf S_{cd}\mathbf M_{ad}\), shows that the
expression is affine in \(\mathbf M\) and in \(\mathbf M\otimes\mathbf M\); the exponential moments are
finite because \(\lVert e^{\mathring\omega}\rVert\leq e^{\lVert\mathring\omega\rVert}\) and Gaussian laws
have exponential moments of every order.
(ii) The differential of \eqref{eq:glink-expected-record} is
\(-(\mu_i-o_e-\mathbf M\mu_j)^\top W_e(d\mathbf M)\mu_j+\Tr(W_e\mathbf M\Sigma_j(d\mathbf M)^\top)=-\langle\mathbf J_e,d\mathbf M\rangle\),
and the chain rule through \Cref{lem:link-exp-derivative} with the adjoint
\(\mathbf L(\mathring\omega_e^\top,\cdot)\) gives the generator gradient.
(iii) The quadratic part
\(\tfrac12\Tr(W_e\mathbf M(\Sigma_j+\mu_j\mu_j^\top)\mathbf M^\top)
=\tfrac12\lVert W_e^{1/2}\mathbf M(\Sigma_j+\mu_j\mu_j^\top)^{1/2}\rVert_{\mathrm F}^2\) is positive definite
in \(\mathbf M\), so the function is strictly convex and coercive, and \(\mathbf J_e=0\) is
\((\mu_i-o_e)\mu_j^\top=\mathbf M(\Sigma_j+\mu_j\mu_j^\top)\), whose unique solution is
\eqref{eq:glink-rank-one}; the trace is \(\mu_j^\top(\Sigma_j+\mu_j\mu_j^\top)^{-1}(\mu_i-o_e)\), and the
Sherman--Morrison identity gives its second form. Every exponential is invertible, so
\(e^{\mathring\omega}\neq\mathbf M^\star\) for \(K\geq2\), which the rest of (iii) assumes. With
\(\operatorname{ad}_{\mathbf B}\mathbf T=[\mathbf B,\mathbf T]\), the identity
\(e^{t\mathbf B}\mathbf Te^{-t\mathbf B}=e^{t\operatorname{ad}_{\mathbf B}}\mathbf T\) writes
\(\mathbf L(\mathbf B,\mathbf T)=\bigl(\int_0^1e^{t\operatorname{ad}_{\mathbf B}}dt\bigr)(\mathbf T)e^{\mathbf B}\), and right
multiplication by \(e^{\mathbf B}\) is invertible. Simultaneous triangularization of the commuting maps
\(\mathbf T\mapsto\mathbf B\mathbf T\) and \(\mathbf T\mapsto\mathbf T\mathbf B\) shows that the eigenvalues of
\(\operatorname{ad}_{\mathbf B}\) are the differences \(\mathring\lambda_a-\mathring\lambda_b\) of eigenvalues of
\(\mathbf B\), so those of \(\int_0^1e^{t\operatorname{ad}_{\mathbf B}}dt\) are
\(\int_0^1e^{t(\mathring\lambda_a-\mathring\lambda_b)}dt\), which vanish exactly when
\(\mathring\lambda_a-\mathring\lambda_b\in2\pi\mathrm i\mathbb Z\setminus\{0\}\); and \(\mathbf B^\top\) has the
eigenvalues of \(\mathbf B\). Where \(\mathbf L(\mathring\omega_e^\top,\cdot)\) is invertible, a zero gradient
forces \(\mathbf J_e(e^{\mathring\omega_e})=0\), hence \(e^{\mathring\omega_e}=\mathbf M^\star\), which is
impossible. If \(\Tr\mathbf M^\star>0\), the rank-one matrix \(\mathbf M^\star\) has the simple eigenvalue
\(\Tr\mathbf M^\star\) and the eigenvalue \(0\) with a \((K-1)\)-dimensional eigenspace, so
\(\mathbf M^\star=\mathbf B\operatorname{diag}(\Tr\mathbf M^\star,0,\ldots,0)\mathbf B^{-1}\) for an invertible
\(\mathbf B\), and
\(\exp\bigl(\mathbf B\operatorname{diag}(\log\Tr\mathbf M^\star,\log\mathrm c_0,\ldots,\log\mathrm c_0)\mathbf B^{-1}\bigr)\to\mathbf M^\star\)
as \(\mathrm c_0\to0\). If \(\Tr\mathbf M^\star=0\) and \(\mathbf M^\star\neq0\), then \((\mathbf M^\star)^2=0\) and
\(\exp\bigl((\log\mathrm c_0)I+\mathrm c_0^{-1}\mathbf M^\star\bigr)=\mathrm c_0I+\mathbf M^\star\to\mathbf M^\star\); if
\(\mathbf M^\star=0\), \(\exp((\log\mathrm c_0)I)\to0\). In each case continuity carries the values to the
infimum, which no exponential attains by strict convexity. If \(\Tr\mathbf M^\star<0\), it is a simple
eigenvalue of \(\mathbf M^\star\); every real matrix near \(\mathbf M^\star\) has a simple real negative
eigenvalue near it, because a simple real eigenvalue of a real matrix stays real under small real
perturbations, and its single Jordan block of size one occurs once, so by
\citet[Theorem~1]{culver1966existence} that matrix has no real logarithm. A neighborhood of
\(\mathbf M^\star\) therefore contains no exponential, and outside it the strictly convex coercive function
stays above a value larger than its minimum.
(iv) At \(\mathbf M=e^{\mathring\omega_0}=-e^{\mathrm a}I_2\),
\(\mu_i-o_e-\mathbf M\mu_j=-2\varsigma^2e^{\mathrm a}(1,0)^\top\), so
\(\mathbf J_e=-2\varsigma^2e^{\mathrm a}\operatorname{diag}(1,0)+\varsigma^2e^{\mathrm a}I_2
=-\varsigma^2e^{\mathrm a}\operatorname{diag}(1,-1)\). The matrix \(\operatorname{diag}(1,-1)\)
anticommutes with \(\mathfrak J\), so
\(e^{\mathrm b\mathfrak J}\operatorname{diag}(1,-1)=\operatorname{diag}(1,-1)e^{-\mathrm b\mathfrak J}\) for real
\(\mathrm b\), and with \(e^{t\mathring\omega_0^\top}=e^{t\mathrm a}e^{-t\pi\mathfrak J}\) the gradient is
\(-\mathbf L(\mathring\omega_0^\top,\mathbf J_e)=\varsigma^2e^{2\mathrm a}\operatorname{diag}(1,-1)\int_0^1e^{(2t-1)\pi\mathfrak J}dt=0\),
since \(\int_{-1}^1\cos(\pi u)du=\int_{-1}^1\sin(\pi u)du=0\). Along the direction \(\operatorname{diag}(1,-1)\),
which anticommutes with \(\mathfrak J\),
\((\pi\mathfrak J+t\operatorname{diag}(1,-1))^2=-(\pi^2-t^2)I_2\), so with \(\mathrm r_t=(\pi^2-t^2)^{1/2}\),
\(e^{\mathring\omega_0+t\operatorname{diag}(1,-1)}
=e^{\mathrm a}\bigl[\cos\mathrm r_tI_2+\mathrm r_t^{-1}\sin\mathrm r_t(\pi\mathfrak J+t\operatorname{diag}(1,-1))\bigr]
=-e^{\mathrm a}I_2+\tfrac{e^{\mathrm a}t^2}{2\pi^2}(\pi\mathfrak J+t\operatorname{diag}(1,-1))+O(t^4)\). The
expected energy is quadratic in \(\mathbf M\) with derivative \(-\mathbf J_e\), and
\(\operatorname{diag}(1,-1)\) is orthogonal to \(\mathfrak J\) and to \(I_2\) for the trace pairing, so
its change is \(\varsigma^2e^{2\mathrm a}t^3/\pi^2+O(t^4)\), negative for small \(t<0\).
\hfill$\square$

Part (iii) is the Gaussian instance of \Cref{cor:link-prior-necessary}(a) with an observed offset: a
single context determines only a rank-one regression, and without the anchor of the generator prior the
record drives the link out of the exponential chart. Part (iv) adds to the stationary generator of
\Cref{cor:link-prior-necessary}(a) that a stationary point at a critical point of the exponential map can
be a saddle of the chart energy. When \(\Tr\mathbf M^\star<0\),
\Cref{prop:glink-chart-infimum-semisimple} excludes attainment of the chart infimum at every chart link
with semisimple negative eigenvalues, which settles \(K=2\) and \(K=3\); attainment at a chart link with a
non-semisimple negative eigenvalue, possible only for \(K\geq4\), is \Cref{open:glink-chart-infimum}.

\propositionheading{Non-attainment of the chart infimum at semisimple negative eigenvalues}{prop:glink-chart-infimum-semisimple}
Let \(K\geq2\), let \(q_i=\mathcal N(\mu_i,\Sigma_i)\) and \(q_j=\mathcal N(\mu_j,\Sigma_j)\) be
nondegenerate, let matter be recognized by \(q_i\otimes q_j\), and, without a link prior, let
\(\mathbf l(\mathbf M)=\E_{q_i\otimes q_j}\mathsf E_e(\cdot;\mathbf M)\) be the expected energy
\eqref{eq:glink-expected-record} of the comparison record with weight \(W_e\succ0\) at an observation
\(o_e\in\R^K\), whose minimizer \(\mathbf M^\star\) of \eqref{eq:glink-rank-one} has \(\Tr\mathbf M^\star<0\).
Call an element of \(\exp(\mathfrak g)\) a chart link, and an eigenvalue semisimple if all its Jordan
blocks have size one. Then no chart link \(\mathbf M\) whose negative eigenvalues are all semisimple
satisfies \(\mathbf l(\mathbf M)=\inf_{\exp(\mathfrak g)}\mathbf l\). For \(K=2\) and \(K=3\) every chart
link has this property, so \(\mathring\omega\mapsto\mathbf l(e^{\mathring\omega})\) does not attain its
infimum over \(\mathfrak g\); for \(K=2\) this is \Cref{cor:link-prior-necessary}(a)(iv) for the matter
law of \((z_i-o_e,z_j)\). \status{ESTABLISHED}

\noindent\emph{Proof.}
By \Cref{prop:glink-sample-record}(iii), \(\mathbf l\) is a strictly convex quadratic function on
\(\R^{K\times K}\) whose gradient vanishes only at \(\mathbf M^\star\), which is singular because its rank is
at most one and \(K\geq2\), while every chart link is invertible, since
\(\det e^{\mathring\omega}=e^{\Tr\mathring\omega}\). By \citet[Theorem~1]{culver1966existence}, a real
\(K\times K\) matrix is a chart link exactly when it is invertible and every Jordan block belonging to a
negative eigenvalue occurs an even number of times. Let a chart link \(\mathbf M\) with semisimple negative
eigenvalues satisfy \(\mathbf l(\mathbf M)=\inf_{\exp(\mathfrak g)}\mathbf l\). It suffices to find a
nonempty open set of real matrices \(\mathbf E\) such that, for each of them, \(\mathbf M+t\mathbf E\) is a
chart link for every sufficiently small \(\lvert t\rvert\). Indeed, for such \(\mathbf E\) minimality gives
\(\mathbf l(\mathbf M+t\mathbf E)\geq\mathbf l(\mathbf M)\) for small \(t\) of either sign, so
\(\langle\nabla\mathbf l(\mathbf M),\mathbf E\rangle=\frac{d}{dt}\mathbf l(\mathbf M+t\mathbf E)\big|_{t=0}=0\);
the linear span of a nonempty open set is \(\R^{K\times K}\), so \(\nabla\mathbf l(\mathbf M)=0\) and
\(\mathbf M=\mathbf M^\star\), which contradicts the invertibility of \(\mathbf M\).

If \(\mathbf M\) has no negative eigenvalue, its spectrum is a compact set at positive distance from
\((-\infty,0]\); by continuity of the roots of the characteristic polynomial every real matrix near
\(\mathbf M\) has no eigenvalue in \((-\infty,0]\) and is a chart link by the criterion, so every
\(\mathbf E\) qualifies. Otherwise let \(\mathrm c'_1,\ldots,\mathrm c'_{\mathrm r}\) be the distinct negative
eigenvalues of \(\mathbf M\). The Jordan blocks of \(\mathrm c'_l\) have size one, so their number is its
multiplicity \(\mathrm n_l\), which the criterion makes even. The real eigenspaces of the \(\mathrm c'_l\) and
the sum of the real generalized eigenspaces of the other eigenvalues decompose \(\R^K\) into invariant
subspaces, so that for a real invertible \(\mathbf S\)
\[
\mathbf S^{-1}\mathbf M\mathbf S=\operatorname{blockdiag}\bigl(\mathrm c'_1I_{\mathrm n_1},\ldots,\mathrm c'_{\mathrm r}I_{\mathrm n_{\mathrm r}},\mathbf M_0\bigr),
\]
with \(\mathbf M_0\) real and without eigenvalues in \((-\infty,0]\). For a real \(\mathbf E\), let
\(\widetilde{\mathbf E}_{ll}\) be the diagonal block of \(\mathbf S^{-1}\mathbf E\mathbf S\) at \(\mathrm c'_l\), the
compression of \(\mathbf E\) to the eigenspace of \(\mathrm c'_l\) along this decomposition, and call
\(\mathbf E\) admissible if every \(\widetilde{\mathbf E}_{ll}\) has distinct non-real eigenvalues. The
admissible matrices form an open set, because each \(\mathbf E\mapsto\widetilde{\mathbf E}_{ll}\) is linear
and, by continuity of the roots of the characteristic polynomial, distinct non-real eigenvalues stay
distinct and non-real under small real perturbations; and a nonempty one, because each \(\mathrm n_l\) is
even and the block-diagonal matrix with the \(2\times2\) blocks
\(\begin{psmallmatrix}0&-\mathrm k\\\mathrm k&0\end{psmallmatrix}\), \(\mathrm k=1,\ldots,\mathrm n_l/2\), has
the distinct eigenvalues \(\pm\mathrm i\mathrm k\).

Fix an admissible \(\mathbf E\), and let \(\mathrm r_0>0\) be smaller than the distance between any two
distinct eigenvalues of \(\mathbf M\) and than the distance from the spectrum of \(\mathbf M_0\) to
\((-\infty,0]\). Fix \(l\), order the blocks so that the one at \(\mathrm c'_l\) comes first, let \(\mathbf B\)
be the complementary diagonal block of \(\mathbf S^{-1}\mathbf M\mathbf S\), of which \(\mathrm c'_l\) is not an
eigenvalue, and split \(\mathbf S^{-1}\mathbf E\mathbf S\) accordingly into
\(\mathbf E_{11}=\widetilde{\mathbf E}_{ll}\), \(\mathbf E_{12}\), \(\mathbf E_{21}\), and \(\mathbf E_{22}\). For
\(\lvert\mathring\lambda-\mathrm c'_l\rvert\leq\mathrm r_0/2\) and small \(\lvert t\rvert\), the matrix
\(\mathbf T(t,\mathring\lambda)=\mathring\lambda I-\mathbf B-t\mathbf E_{22}\) is invertible, and the Schur
complement gives
\[
\det(\mathring\lambda I-\mathbf M-t\mathbf E)
=\det\mathbf T(t,\mathring\lambda)
\det\bigl((\mathring\lambda-\mathrm c'_l)I-t\mathbf E_{11}-t^2\mathbf E_{12}\mathbf T(t,\mathring\lambda)^{-1}\mathbf E_{21}\bigr).
\]
With \(\mathring\lambda=\mathrm c'_l+t\mathring\nu\) the second factor is \(t^{\mathrm n_l}\mathrm f(t,\mathring\nu)\),
where
\(\mathrm f(t,\mathring\nu)=\det\bigl(\mathring\nu I-\mathbf E_{11}-t\mathbf E_{12}\mathbf T(t,\mathrm c'_l+t\mathring\nu)^{-1}\mathbf E_{21}\bigr)\)
is real-analytic in \(t\) and holomorphic in \(\mathring\nu\) near \(t=0\), and \(\mathrm f(0,\cdot)\) is the
characteristic polynomial of \(\mathbf E_{11}\), whose roots \(\mathring\nu_1,\ldots,\mathring\nu_{\mathrm n_l}\) are
simple and non-real. The implicit function theorem at these simple roots gives continuous
\(\mathring\nu_k(t)\) for small \(\lvert t\rvert\), with \(\mathrm f(t,\mathring\nu_k(t))=0\) and
\(\mathring\nu_k(0)=\mathring\nu_k\), which stay distinct and non-real. For small \(t\neq0\) the numbers
\(\mathrm c'_l+t\mathring\nu_k(t)\) are therefore \(\mathrm n_l\) distinct eigenvalues of \(\mathbf M+t\mathbf E\)
within \(\mathrm r_0/2\) of \(\mathrm c'_l\), and they are non-real, with imaginary parts
\(t\operatorname{Im}\mathring\nu_k(t)\neq0\). By continuity of the roots of the characteristic polynomial,
for small \(\lvert t\rvert\) exactly \(\mathrm n_l\) eigenvalues of \(\mathbf M+t\mathbf E\), counted with multiplicity,
lie within \(\mathrm r_0/2\) of \(\mathrm c'_l\), so these are all of them. The same continuity places every
eigenvalue of \(\mathbf M+t\mathbf E\) within \(\mathrm r_0/2\) of the spectrum of \(\mathbf M\) for small
\(\lvert t\rvert\), and those not near any \(\mathrm c'_l\) lie near the spectrum of \(\mathbf M_0\), off
\((-\infty,0]\). Hence, for every small \(t\neq0\) of either sign, \(\mathbf M+t\mathbf E\) has no eigenvalue in
\((-\infty,0]\) and is a chart link by the criterion.

For \(K\leq3\), a negative eigenvalue of a chart link has even multiplicity by the criterion, hence
multiplicity two, and a single Jordan block of size two would occur once, so it has two blocks of size
one and is semisimple; for \(K=3\) the remaining eigenvalue is real and, as \(\det e^{\mathring\omega}>0\),
positive. The infimum of \(\mathring\omega\mapsto\mathbf l(e^{\mathring\omega})\) over \(\mathfrak g\) is
\(\inf_{\exp(\mathfrak g)}\mathbf l\), attained at \(\mathring\omega\) exactly when \(e^{\mathring\omega}\)
attains it. For \(K=2\), \(\mathbf l\) is, up to the constant \(\tfrac12\log\det(2\pi W_e^{-1})\), the
expected record energy of \Cref{cor:link-prior-necessary}(a) for the matter law of \((z_i-o_e,z_j)\), a
product law with \(\Sigma_j\succ0\), whose condition \(\mu_j^\top\Sigma_j^{-1}(\mu_i-o_e)<0\) is
\(\Tr\mathbf M^\star<0\) by \eqref{eq:glink-rank-one}. \hfill$\square$

\section{Expected loop defects under product recognition}
\label{sec:glink-loops}

This section computes expected loop defects when the generators are recognized by a product Gaussian
law.

\propositionheading{Expected loop defects under product Gaussian recognition}{prop:glink-expected-loop}
Let the generators be recognized by the product Gaussian law
\(R=\bigotimes_{e\in E_+}\mathcal N(\mathring\mu_e,\mathring\Lambda_e^{-1})\), with each precision
\(\mathring\Lambda_e\) a positive-definite self-adjoint operator on
\((\mathfrak g,\langle\cdot,\cdot\rangle_{\mathsf G})\), as in \Cref{prop:link-gaussian-recognition-vfe}, and
for a copy \(e\notin E_+\) write \(\mathring\mu_e=-\mathring\mu_{\bar e}\) and
\(\mathring\Lambda_e=\mathring\Lambda_{\bar e}\), so that its generator \(\mathring\omega_e=-\mathring\omega_{\bar e}\)
has the law \(\mathcal N(\mathring\mu_e,\mathring\Lambda_e^{-1})\) under \(R\). Let \(p=(e_0,\ldots,e_{r-1})\) be
a closed walk whose copies belong to pairwise distinct reversal pairs. (i) The expected reference
holonomy and its second moment are the ordered products
\begin{equation}
\E_R\mathring\Theta_p=\prod_{a=0}^{r-1}\E_Re^{\mathring\omega_{e_a}},
\qquad
\E_R\bigl[\mathring\Theta_p\otimes\mathring\Theta_p\bigr]
=\prod_{a=0}^{r-1}\E_R\bigl[e^{\mathring\omega_{e_a}}\otimes e^{\mathring\omega_{e_a}}\bigr],
\label{eq:glink-loop-moments}
\end{equation}
with factors in the order of the walk. (ii) With \(\mathsf G=\mathsf G^{\mathrm{fib}}\) and the vectors
\(\mathbf G^\pm\in\R^K\otimes\R^K\) of coordinates \(\mathbf G^+_{(a,b)}=\mathsf G_{ab}\) and
\(\mathbf G^-_{(a,b)}=(\mathsf G^{-1})_{ab}\), the expected loop defect is the finite expression
\begin{equation}
\E_R\mathfrak Q_p
=\tfrac12(\mathbf G^+)^\top\Bigl(\prod_{a=0}^{r-1}\E_R\bigl[e^{\mathring\omega_{e_a}}\otimes e^{\mathring\omega_{e_a}}\bigr]\Bigr)\mathbf G^-
-\Tr\Bigl(\prod_{a=0}^{r-1}\E_Re^{\mathring\omega_{e_a}}\Bigr)+\tfrac K2
\label{eq:glink-expected-defect}
\end{equation}
in the per-copy first and second moments of exponentials of Gaussian matrices.
(iii) Write \(\mathring\omega_{e_a}=\mathring\mu_{e_a}+\mathring\zeta_a\), with \(\mathring\mu_{e_a}\) the mean of
the generator of \(e_a\), and let \(\mathring\Theta_p(\mathring\mu)\),
\(\mathring\Theta^{<a}_p=\prod_{b<a}e^{\mathring\mu_{e_b}}\), and
\(\mathring\Theta^{>a}_p=\prod_{b>a}e^{\mathring\mu_{e_b}}\) be evaluated at the means. As the covariances
\(\mathring\Lambda_e^{-1}\) of the copies of \(p\) tend to zero with the means held,
\begin{equation}
\begin{aligned}
\E_R\mathfrak Q_p={}&\mathfrak Q_p(\mathring\mu)
+\tfrac12\sum_{a=0}^{r-1}\E\Bigl[\bigl\lVert\mathring\Theta^{<a}_p\mathbf L(\mathring\mu_{e_a},\mathring\zeta_a)\mathring\Theta^{>a}_p\bigr\rVert_{\mathsf G}^2\\
&+\bigl\langle\mathring\Theta_p(\mathring\mu)-I,\mathring\Theta^{<a}_p\mathbf L^{(2)}(\mathring\mu_{e_a};\mathring\zeta_a)\mathring\Theta^{>a}_p\bigr\rangle_{\mathsf G}\Bigr]
+O\Bigl(\max_a\lVert\mathring\Lambda_{e_a}^{-1}\rVert^2\Bigr),
\end{aligned}
\label{eq:glink-defect-first-order}
\end{equation}
and at flat means, \(\mathring\mu_{e_a}=0\) for every \(a\),
\(\E_R\mathfrak Q_p=\tfrac12\sum_a\Tr\mathring\Lambda_{e_a}^{-1}+O(\max_a\lVert\mathring\Lambda_{e_a}^{-1}\rVert^2)\),
where \(\Tr\) is the trace of an operator on \((\mathfrak g,\langle\cdot,\cdot\rangle_{\mathsf G})\).
(iv) If a walk traverses a reversal pair twice, (i) fails in general: for \(K=1\) and the walk
\((e,\bar e)\), \(\mathring\Theta_p=1\) while
\(\E_Re^{\mathring\omega_e}\E_Re^{-\mathring\omega_e}=e^{\mathring\Lambda_e^{-1}}\neq1\).
\status{ESTABLISHED}

\noindent\emph{Proof.}
(i) Distinct reversal pairs carry independent generators under \(R\), so the factors
\(e^{\mathring\omega_{e_a}}\) are independent random matrices. They are integrable to every order,
because \(\lVert e^{\mathbf B}\rVert\leq e^{\lVert\mathbf B\rVert}\) and Gaussian laws have exponential moments of
every order. For independent integrable random matrices \(\E[\mathbf S\mathbf T]=(\E\mathbf S)(\E\mathbf T)\)
entrywise, which gives the first product, and the mixed-product rule
\((\mathbf B_0\cdots\mathbf B_{r-1})\otimes(\mathbf B_0\cdots\mathbf B_{r-1})=(\mathbf B_0\otimes\mathbf B_0)\cdots(\mathbf B_{r-1}\otimes\mathbf B_{r-1})\)
reduces the second to the first. (ii) In the coordinates of \Cref{sec:glink-setting},
\((\mathbf G^+)^\top(\mathbf B\otimes\mathbf B)\mathbf G^-=\sum_{a,b,c,d}\mathsf G_{ab}\mathbf B_{ac}\mathbf B_{bd}(\mathsf G^{-1})_{cd}
=\Tr(\mathsf G^{-1}\mathbf B^\top\mathsf G\mathbf B)=\langle\mathbf B,\mathbf B\rangle_{\mathsf G}\). Expanding
\(\tfrac12\langle\mathring\Theta_p-I,\mathring\Theta_p-I\rangle_{\mathsf G}\) with
\(\langle I,\mathbf B\rangle_{\mathsf G}=\Tr\mathbf B\) and \(\langle I,I\rangle_{\mathsf G}=K\), and taking
expectations with (i), gives \eqref{eq:glink-expected-defect}.
(iii) At the generators \(\mathring\mu_{e_a}+t\mathring\zeta_a\), the defect \(\mathfrak Q_p\) is a polynomial in
the entries of the factors \(e^{\mathring\mu_{e_a}+t\mathring\zeta_a}\), and the power series of the
exponential bounds the \(k\)-th derivative of \(t\mapsto e^{\mathbf B+t\mathbf T}\) by
\(\lVert\mathbf T\rVert^ke^{\lVert\mathbf B\rVert+\lvert t\rvert\lVert\mathbf T\rVert}\); hence, with
\(\mathring\zeta=(\mathring\zeta_0,\ldots,\mathring\zeta_{r-1})\), the fourth \(t\)-derivative of the defect is
bounded by \(\mathrm c_1e^{\mathrm c_1\lVert\mathring\zeta\rVert}\lVert\mathring\zeta\rVert^4\) for
\(\lvert t\rvert\leq1\), with \(\mathrm c_1\) depending on the means. Taylor's formula in \(t\) to third order
with integral remainder at \(t=1\), the vanishing of first and third moments of a centered Gaussian, and
the bound
\(\E[e^{\mathrm c_1\lVert\mathring\zeta\rVert}\lVert\mathring\zeta\rVert^4]=O(\max_a\lVert\mathring\Lambda_{e_a}^{-1}\rVert^2)\)
for bounded covariances give \(\E_R\mathfrak Q_p=\mathfrak Q_p(\mathring\mu)\) plus one half the expected
second \(t\)-derivative at \(t=0\), up to \(O(\max_a\lVert\mathring\Lambda_{e_a}^{-1}\rVert^2)\). In that second
derivative the terms mixing \(\mathring\zeta_a\) and \(\mathring\zeta_b\), \(a\neq b\), have zero expectation
because the two are independent and centered. Along \(\mathring\zeta_a\) alone, the first and second
derivatives of \(\mathring\Theta_p\) are \(\mathring\Theta^{<a}_p\mathbf L(\mathring\mu_{e_a},\mathring\zeta_a)\mathring\Theta^{>a}_p\)
and \(\mathring\Theta^{<a}_p\mathbf L^{(2)}(\mathring\mu_{e_a};\mathring\zeta_a)\mathring\Theta^{>a}_p\), and the second
derivative of \(\mathfrak Q_p\) is the squared \(\mathsf G\)-norm of the first plus the pairing of
\(\mathring\Theta_p-I\) with the second, which gives \eqref{eq:glink-defect-first-order}. At flat means
\(\mathring\Theta^{<a}_p=\mathring\Theta^{>a}_p=I\), \(\mathbf L(0,\mathring\zeta_a)=\mathring\zeta_a\), and
\(\mathring\Theta_p(0)=I\), so the bracket is \(\E\lVert\mathring\zeta_a\rVert_{\mathsf G}^2=\Tr\mathring\Lambda_{e_a}^{-1}\).
(iv) For \(K=1\) and \(\mathring\omega_e\sim\mathcal N(\mathring\mu_e,\mathring\Lambda_e^{-1})\),
\(\E e^{\pm\mathring\omega_e}=e^{\pm\mathring\mu_e+\mathring\Lambda_e^{-1}/2}\). \hfill$\square$

At flat means the recognition spread alone produces curvature energy: each copy of the walk
contributes half its generator variance, measured in the fiber metric. For a repeated pair the
expectation involves joint moments of one Gaussian generator, such as
\(\E[e^{\mathring\omega_e}\mathbf Te^{\mathring\omega_e}]\), which the same Taylor argument expands to first
order but which do not factor.

\section{The action induced by undirected Gaussian matter}
\label{sec:glink-induced}

On a finite multigraph with agents \(V=\{1,\ldots,N\}\), let every copy \(e:j\to i\) in \(E_+\) carry a
weight \(W_e\succ0\), every agent a self term \(A_i\succeq0\), and let \(\mathsf h\in\R^{NK}\) be an
information vector, all in reference coordinates. The matter energy with latent links is
\begin{equation}
\mathsf E(z;\mathring\omega)
=\tfrac12\sum_{i\in V}z_i^\top A_iz_i
+\tfrac12\sum_{e\in E_+}\bigl(z_i-e^{\mathring\omega_e}z_j\bigr)^\top W_e\bigl(z_i-e^{\mathring\omega_e}z_j\bigr)
-\mathsf h^\top z,
\label{eq:glink-matter-energy}
\end{equation}
with precision \(\Lambda(\mathring\omega)=A_\oplus+L(\mathring\omega)\), where
\(A_\oplus=\operatorname{blockdiag}(A_1,\ldots,A_N)\) and \(L(\mathring\omega)\) is the connection
Laplacian of \Cref{def:cg-connection-laplacian} with represented links \(e^{\mathring\omega_e}\); a loop
copy contributes \((I-e^{\mathring\omega_e})^\top W_e(I-e^{\mathring\omega_e})\) to its agent's block.
At \(\mathring\omega=0\) the precision is the member of the interaction family
\eqref{eq:gauss-interaction-family} with self terms \(A_i\) and weights summed over parallel copies,
the form adopted in \Cref{hyp:gauss-global-interaction}; for \(\mathring\omega\neq0\) it has the
nonflat form of \Cref{def:cg-nonflat-series-class}. The generator family carries a base law
\(e^{-\mathscr S}\nu^{\mathscr T}\) with \(0<Z_{\mathscr S}<\infty\) as in
\Cref{prop:elbo-random-link-vfe}, for instance the generator prior of \Cref{def:link-generator-prior},
and the joint law of generators and matter is proportional to
\(e^{-\mathscr S(\mathring\omega)}\nu^{\mathscr T}(d\mathring\omega)e^{-\mathsf E(z;\mathring\omega)}dz\),
normalized as a whole, the undirected model of \Cref{def:gen-undirected-model}.

\propositionheading{Induced action of undirected Gaussian matter}{prop:glink-induced-action}
(i) If every \(A_i\succ0\), then \(\Lambda(\mathring\omega)\succeq A_\oplus\succ0\) for every
\(\mathring\omega\), the joint normalizer \(\mathfrak Z\), as in \Cref{prop:link-induced-action}, is finite
and positive, and the generator marginal is
\(\mathfrak Z^{-1}e^{-\mathscr S(\mathring\omega)-\mathscr S^{\mathrm{ind}}(\mathring\omega)}\nu^{\mathscr T}(d\mathring\omega)\) with
\begin{equation}
\mathscr S^{\mathrm{ind}}(\mathring\omega)=-\mathsf A\bigl(\mathsf h,\Lambda(\mathring\omega)\bigr)
=\tfrac12\log\det\Lambda(\mathring\omega)-\tfrac12\mathsf h^\top\Lambda(\mathring\omega)^{-1}\mathsf h-\tfrac{NK}2\log2\pi
\geq-\mathsf A(\mathsf h,A_\oplus).
\label{eq:glink-induced-action}
\end{equation}
(ii) If every \(A_i\succ0\) and \(E_+\neq\varnothing\), then \(e^{-\mathscr S^{\mathrm{ind}}}\) is not
integrable for Lebesgue measure on \(\mathscr T\); the induced action does not by itself define a law
of the generators.
(iii) If the multigraph is a forest (no loop copy, no parallel copies, no cycle), choose a root in
every component, let \(\mathring\Theta_{I\leftarrow i}\) be the product of the twists along the path from \(i\)
to the root of its component \(I\), the tree transport of \Cref{def:rg-rooted-tree-data} in reference
coordinates, and put \(\mathring\Theta_\oplus=\operatorname{blockdiag}(\mathring\Theta_{I\leftarrow i})\). Then
\(\Lambda(\mathring\omega)=\mathring\Theta_\oplus^\top\Lambda^{\mathrm{root}}(\mathring\omega)\mathring\Theta_\oplus\), where
\(\Lambda^{\mathrm{root}}(\mathring\omega)\) belongs to the interaction family with self terms
\(\mathring\Theta_{I\leftarrow i}^{-\top}A_i\mathring\Theta_{I\leftarrow i}^{-1}\) and weights
\(\mathring\Theta_{I\leftarrow i}^{-\top}W_e\mathring\Theta_{I\leftarrow i}^{-1}\) on \(e:j\to i\); no holonomy
enters \(\mathscr S^{\mathrm{ind}}\). If moreover \(A_i=a_i\mathsf G\), \(W_e=w_e\mathsf G\) with scalars
\(a_i,w_e>0\), \(\mathsf h=0\), and every link is an isometry,
\(e^{\mathring\omega_e^\top}\mathsf Ge^{\mathring\omega_e}=\mathsf G\), then \(\mathscr S^{\mathrm{ind}}\)
takes one value on all such configurations.
(iv) For two agents joined by one copy \(e:j\to i\) and \(\mathsf h=0\),
\(\det\Lambda(\mathring\omega)=\det(A_i+W_e)\det(A_j+e^{\mathring\omega_e^\top}W_e^\sharp e^{\mathring\omega_e})\)
with \(W_e^\sharp=W_e(A_i+W_e)^{-1}A_i\). For \(\mathsf G=I_K\), \(A_i=a_iI\) with \(a_i>0\), \(A_j=a_jI\),
\(W_e=w_eI\), \(w^\sharp=a_iw_e/(a_i+w_e)\), and \(\mathring\omega_e=t\mathring\upsilon\) with \(\mathring\upsilon\)
symmetric,
\begin{equation}
\frac{d}{dt}\mathscr S^{\mathrm{ind}}\Big|_{t=0}=\frac{w^\sharp\Tr\mathring\upsilon}{a_j+w^\sharp},
\qquad
\frac{d^2}{dt^2}\mathscr S^{\mathrm{ind}}\Big|_{t=0}=\frac{2w^\sharp a_j\lVert\mathring\upsilon\rVert_{\mathrm F}^2}{(a_j+w^\sharp)^2},
\label{eq:glink-edge-mass}
\end{equation}
while \(\mathscr S^{\mathrm{ind}}\) is constant along skew-symmetric \(\mathring\upsilon\): noncompact generator
directions receive edge-local terms without any loop, and along the trace direction the action
decreases as the link contracts.
(v) At fixed links the induced action takes the values of earlier chapters. With every coordinate
hidden it is the scalar \(\mathfrak c^{\mathrm{eff}}-\mathfrak c\) of \eqref{eq:cg-kron-complete-schur}
at the links \(e^{\mathring\omega}\). For an unanchored cycle of \(m\geq2\) distinct agents, with copies
\(e_a:i_{a+1}\to i_a\), \(i_m=i_0\), and every \(A_i=0\),
\begin{equation}
\det\Lambda(\mathring\omega)=\det\bigl(I-\mathring\Theta_p\bigr)^2\prod_{a=0}^{m-1}\det W_{e_a},
\qquad p=(e_0,\ldots,e_{m-1}),
\label{eq:glink-cycle-determinant}
\end{equation}
which for \(m=2\) and \(W_e=\Sigma_e^{-1}\) is \eqref{eq:obs-holonomy-det}. For \(K=1\), the reciprocal
pair of \Cref{prop:obs-normalizer-link-dependence} with unit weights, anchors \(a_0>0\), and the flat
slice \(\mathring\omega_{e'}=-\mathring\omega_e\) has
\(\mathscr S^{\mathrm{ind}}=\tfrac12\log(a_0^2+4a_0\cosh^2\mathring\omega_e)-\log2\pi\) at
\(\mathsf h=0\), which is \eqref{eq:obs-det-witness} with \(\mathsf g=e^{\mathring\omega_e}\).
\status{ESTABLISHED}

\noindent\emph{Proof.}
(i) Every edge term of \(L(\mathring\omega)\) is positive semidefinite, so
\(\Lambda(\mathring\omega)\succeq A_\oplus\succ0\). By \Cref{prop:gauss-log-normalizer-moments},
\(\int e^{-\mathsf E(z;\mathring\omega)}dz=\exp\mathsf A(\mathsf h,\Lambda(\mathring\omega))\). The Loewner
inequality gives \(\log\det\Lambda(\mathring\omega)\geq\log\det A_\oplus\) and
\(\Lambda(\mathring\omega)^{-1}\preceq A_\oplus^{-1}\), so
\(\mathsf A(\mathsf h,\Lambda(\mathring\omega))\leq\mathsf A(\mathsf h,A_\oplus)\), which is the lower bound.
The map \(\mathring\omega\mapsto\mathsf A(\mathsf h,\Lambda(\mathring\omega))\) is continuous, so by
Tonelli's theorem
\(\mathfrak Z=\int e^{-\mathscr S}e^{\mathsf A(\mathsf h,\Lambda)}d\nu^{\mathscr T}\in(0,Z_{\mathscr S}e^{\mathsf A(\mathsf h,A_\oplus)}]\),
and integrating out \(z\) gives the generator marginal.
(ii) Fix \(e_\star\in E_+\) and consider the generator families with
\(\mathring\omega_{e_\star}=tI+\mathring\xi\), \(t\leq0\), \(\lVert\mathring\xi\rVert\leq1\), and
\(\lVert\mathring\omega_e\rVert\leq1\) for \(e\neq e_\star\), in the operator norm. On this set every link lies in
a compact set of matrices: \(e^{\mathring\omega_e}\) in the image of the unit ball, and
\(e^{tI+\mathring\xi}=e^te^{\mathring\xi}\) among the multiples \(\mathrm a\mathbf B\) with \(\mathrm a\in[0,1]\) and
\(\mathbf B\) in that image. The precision is a polynomial in the links and is \(\succeq A_\oplus\) at every
matrix value of the links, singular ones included, so \(\tfrac12\log\det\Lambda\) is bounded above on the
set, and \(\mathscr S^{\mathrm{ind}}\leq\tfrac12\log\det\Lambda-\tfrac{NK}2\log2\pi\). The set contains the
disjoint sets on which \(\lVert\mathring\omega_{e_\star}+kI\rVert\leq\tfrac13\) and
\(\lVert\mathring\omega_e\rVert\leq1\) for \(e\neq e_\star\), \(k=1,2,\ldots\), of equal positive Lebesgue
measure, so it has infinite measure, and \(e^{-\mathscr S^{\mathrm{ind}}}\) is bounded below by a positive
constant on it.
(iii) The transports exist and are unique given the roots, because each copy of a forest joins a vertex
to its parent and the relation \(\mathring\Theta_{I\leftarrow i}e^{\mathring\omega_e}=\mathring\Theta_{I\leftarrow j}\) on
\(e:j\to i\) determines the child's transport from the parent's in either orientation. With
\(z=\mathring\Theta_\oplus^{-1}z^{\mathrm{root}}\), every residual becomes
\(z_i-e^{\mathring\omega_e}z_j=\mathring\Theta_{I\leftarrow i}^{-1}z^{\mathrm{root}}_i-e^{\mathring\omega_e}\mathring\Theta_{I\leftarrow j}^{-1}z^{\mathrm{root}}_j
=\mathring\Theta_{I\leftarrow i}^{-1}(z^{\mathrm{root}}_i-z^{\mathrm{root}}_j)\), so
\(\mathring\Theta_\oplus^{-\top}\Lambda(\mathring\omega)\mathring\Theta_\oplus^{-1}\) is the interaction-family precision
\(\Lambda^{\mathrm{root}}\) with the stated self terms and weights. Isometric links make each
\(\mathring\Theta_{I\leftarrow i}\) an isometry of \(\mathsf G\), so the self terms and weights of \(\Lambda^{\mathrm{root}}\)
are \(A_i\) and \(W_e\), and \((\det\mathring\Theta_{I\leftarrow i})^2=1\) because
\(\det(\mathring\Theta_{I\leftarrow i}^\top\mathsf G\mathring\Theta_{I\leftarrow i})=\det\mathsf G\); hence
\(\det\Lambda(\mathring\omega)=\det\Lambda^{\mathrm{root}}(0)\).
(iv) With \(\mathring\Theta_e=e^{\mathring\omega_e}\), the blocks are \(\Lambda_{ii}=A_i+W_e\),
\(\Lambda_{ij}=-W_e\mathring\Theta_e\), and \(\Lambda_{jj}=A_j+\mathring\Theta_e^\top W_e\mathring\Theta_e\), and the
Schur complement of \(\Lambda_{ii}\) is \(A_j+\mathring\Theta_e^\top(W_e-W_e(A_i+W_e)^{-1}W_e)\mathring\Theta_e\),
where \(W_e-W_e(A_i+W_e)^{-1}W_e=W_e(A_i+W_e)^{-1}(A_i+W_e-W_e)=W_e(A_i+W_e)^{-1}A_i\), with
\(A_i+W_e\succ0\) because \(W_e\succ0\); \Cref{lem:gauss-schur}(iv) gives the determinant. In the
isotropic case \(\mathscr S^{\mathrm{ind}}=\tfrac12\log\det\mathbf B_t\) plus a constant, with \(\mathbf B_t=a_jI+w^\sharp e^{2t\mathring\upsilon}\), which commutes with \(\mathring\upsilon\); the
derivatives \(\frac{d}{dt}\log\det\mathbf B_t=\Tr(\mathbf B_t^{-1}\mathbf B_t')\) and
\(\frac{d^2}{dt^2}\log\det\mathbf B_t=\Tr(\mathbf B_t^{-1}\mathbf B_t''-\mathbf B_t^{-1}\mathbf B_t'\mathbf B_t^{-1}\mathbf B_t')\) at
\(t=0\), where \(\mathbf B_0=(a_j+w^\sharp)I\), \(\mathbf B_0'=2w^\sharp\mathring\upsilon\), and
\(\mathbf B_0''=4w^\sharp\mathring\upsilon^2\), give \eqref{eq:glink-edge-mass}. For skew \(\mathring\upsilon\),
\(e^{t\mathring\upsilon^\top}e^{t\mathring\upsilon}=I\).
(v) With the pivot \(\Lambda\), \(n_{\mathrm h}=NK\), and no retained coordinate,
\eqref{eq:cg-kron-complete-schur} gives
\(\mathfrak c^{\mathrm{eff}}-\mathfrak c=\tfrac12\log\det\Lambda-\tfrac{NK}2\log2\pi-\tfrac12\mathsf h^\top\Lambda^{-1}\mathsf h\).
For the cycle, \(z^\top\Lambda z=(\mathbf Tz)^\top W_\oplus(\mathbf Tz)\) with
\(W_\oplus=\operatorname{blockdiag}(W_{e_a})\) and the square block matrix \(\mathbf T\) whose block row \(a\)
has \(I\) in column \(a\) and \(-\mathring\Theta_{e_a}\) in column \(a+1\) modulo \(m\), so
\(\det\Lambda=(\det\mathbf T)^2\det W_\oplus\). Deleting block row and column \(0\) of \(\mathbf T\) leaves a
block upper bidiagonal matrix with identity diagonal blocks, of determinant one, whose inverse has the
block \(\mathring\Theta_{e_1}\cdots\mathring\Theta_{e_{m-2}}\) in its upper right corner (the identity when
\(m=2\)); the Schur complement of that matrix in \(\mathbf T\) is
\(I-\mathring\Theta_{e_0}\mathring\Theta_{e_1}\cdots\mathring\Theta_{e_{m-1}}=I-\mathring\Theta_p\), which gives
\eqref{eq:glink-cycle-determinant}. The last statement is substitution into \eqref{eq:obs-det-witness},
since \((e^{\mathring\omega}+e^{-\mathring\omega})^2=4\cosh^2\mathring\omega\). \hfill$\square$

Part (ii) is the Gaussian counterpart of the requirement that the generator prior be normalized: the
induced action is bounded below but stays bounded on a set of infinite volume toward contracting links,
so it confines no generator direction by itself, and the base law \(\nu^{\mathscr T}\) must. Parts (iii)
and (iv) separate the two kinds of generator direction. Compact directions enter
\(\mathscr S^{\mathrm{ind}}\) only through loops, while noncompact directions change the diagonal blocks
of \(\Lambda\) on a single copy; along the trace direction the induced action decreases as links
contract, toward the singular boundary that \Cref{prop:glink-sample-record}(iii) meets without a prior.

\corollaryheading{An unanchored cycle admits no joint law with latent links}{cor:glink-unanchored-cycle}
Let \(m\geq2\) distinct agents form a cycle \(p=(e_0,\ldots,e_{m-1})\) of copies from distinct
reversal pairs, with weights \(W_{e_a}\succ0\), every self term zero, any information vector
\(\mathsf h\), and latent generators of the cycle copies. Let \(\nu\) be a finite measure on
\(\mathfrak g^m\) whose density is bounded below by a positive constant near some point at which
\(\det(I-\mathring\Theta_p)\) vanishes with nonzero differential; for every \(K\geq1\) the point with
\(\mathring\omega_{e_1}=\operatorname{diag}(0,\log2,\ldots,\log K)\) and every other generator zero is
such a point, and every Gaussian law with positive-definite covariance qualifies. Then
\(\int\nu(d\mathring\omega)\int e^{-\mathsf E(z;\mathring\omega)}dz=+\infty\), so no law proportional to
\(\nu(d\mathring\omega)e^{-\mathsf E(z;\mathring\omega)}dz\) exists, although for every generator family
off the Lebesgue-null set \(\{\det(I-\mathring\Theta_p)=0\}\) the matter law at fixed links exists. With
positive-definite self terms at every agent the joint law exists for every finite base measure
(\Cref{prop:glink-induced-action}(i)). \status{ESTABLISHED}

\noindent\emph{Proof.}
Write \(\mathrm D_p(\mathring\omega)=\det(I-\mathring\Theta_p)\). Off its zero set,
\eqref{eq:glink-cycle-determinant} gives \(\Lambda(\mathring\omega)\succ0\), and since
\(\mathsf h^\top\Lambda^{-1}\mathsf h\geq0\),
\(\int e^{-\mathsf E}dz=e^{\mathsf A(\mathsf h,\Lambda)}\geq(2\pi)^{mK/2}(\det\Lambda)^{-1/2}
=(2\pi)^{mK/2}\prod_a(\det W_{e_a})^{-1/2}\lvert\mathrm D_p(\mathring\omega)\rvert^{-1}\); on the zero set
the inner integral is \(+\infty\). The function \(\mathrm D_p\) is real-analytic and not identically
zero, so its zero set is Lebesgue-null. At a point where \(\mathrm D_p=0\) and
\(d\mathrm D_p\neq0\), the submersion theorem gives local coordinates on a neighborhood whose first
coordinate is \(\mathrm D_p\), with Jacobian bounded above and below on a smaller neighborhood, and the
density of \(\nu\) is bounded below there, so the integral of \(\lvert\mathrm D_p\rvert^{-1}\) is bounded
below by a positive multiple of \(\int_{-\mathrm a}^{\mathrm a}\lvert\mathrm u\rvert^{-1}d\mathrm u=+\infty\) for some
\(\mathrm a>0\). At the stated point, along \(\mathring\omega_{e_0}=tI\),
\(\mathring\Theta_p=e^t\operatorname{diag}(1,2,\ldots,K)\) and \(\mathrm D_p=\prod_{k=1}^K(1-ke^t)\), which
vanishes at \(t=0\) with derivative \(-\prod_{k=2}^K(1-k)\neq0\). A Gaussian law with positive-definite
covariance has a continuous positive density. \hfill$\square$

At fixed links, \Cref{prop:obs-reciprocal-pair-kernel,cor:obs-holonomy-kernel-shrinkage} show that
holonomy without eigenvalue one restores definiteness of the unanchored pair, so its law exists off a
null set of links; the corollary shows that this does not survive integration over the links, because
the partition function diverges like \(\lvert\det(I-\mathring\Theta_p)\rvert^{-1}\) across a
hypersurface. For \(K=1\) more cycles restore integrability (\Cref{prop:glink-unanchored-scalar}); for
\(K\geq2\) whether they do is \Cref{open:glink-unanchored-normalization}.

\propositionheading{Joint normalization of unanchored scalar matter}{prop:glink-unanchored-scalar}
Let \(K=1\), so that generators and links are real numbers, and in \eqref{eq:glink-matter-energy} let the
multigraph be connected, every weight \(W_e=w_e\) positive, every self term zero, and the information
vector zero. Let the generators carry a Gaussian base law \(\nu\) on \(\mathscr T=\R^{E_+}\) with
positive-definite covariance and no action, for instance a product of nondegenerate Gaussian laws of the
generators. With the value \(+\infty\) for \(\det\Lambda(\mathring\omega)^{-1/2}\) where \(\Lambda(\mathring\omega)\)
is singular, the joint normalizer of \Cref{prop:link-induced-action} is
\begin{equation}
\mathfrak Z=\int\nu(d\mathring\omega)\int e^{-\mathsf E(z;\mathring\omega)}dz
=(2\pi)^{N/2}\int\nu(d\mathring\omega)\det\Lambda(\mathring\omega)^{-1/2}.
\label{eq:glink-unanchored-normalizer}
\end{equation}
(i) If the cycle rank \(\lvert E_+\rvert-N+1\) of the multigraph, the number of its copies in \(E_+\)
outside a spanning tree, is at least two, then \(\Lambda(\mathring\omega)\succ0\) off a Lebesgue-null set
and \(0<\mathfrak Z<\infty\), so generators and matter have the joint law
\(\mathfrak Z^{-1}\nu(d\mathring\omega)e^{-\mathsf E(z;\mathring\omega)}dz\).
(ii) If the cycle rank is one, then \(\mathfrak Z=+\infty\) and no law proportional to
\(\nu(d\mathring\omega)e^{-\mathsf E(z;\mathring\omega)}dz\) exists; for a cycle of distinct agents this is
\Cref{cor:glink-unanchored-cycle} at \(K=1\). \status{ESTABLISHED}

\noindent\emph{Proof.}
Where \(\Lambda(\mathring\omega)\succ0\), \Cref{prop:gauss-log-normalizer-moments} gives
\(\int e^{-\mathsf E}dz=e^{\mathsf A(0,\Lambda)}=(2\pi)^{N/2}\det\Lambda^{-1/2}\); where \(\Lambda\) is singular,
\(\mathsf E\) is constant along a kernel vector and the inner integral is \(+\infty\); Tonelli's theorem
gives \eqref{eq:glink-unanchored-normalizer}. Let \(\mathbf T(\mathring\omega)\) be the
\(\lvert E_+\rvert\times N\) matrix whose row at a copy \(e:j\to i\) has the entry \(1\) in column \(i\), the
entry \(-e^{\mathring\omega_e}\) in column \(j\), and zeros elsewhere, the row of a loop copy having the
single entry \(1-e^{\mathring\omega_e}\). Then
\(2\mathsf E(z;\mathring\omega)=\sum_{e\in E_+}w_e((\mathbf Tz)_e)^2\), so
\(\Lambda=\mathbf T^\top\operatorname{diag}(w_e)\mathbf T\), and the Cauchy--Binet formula gives
\[
\det\Lambda(\mathring\omega)=\sum_{E'}\Bigl(\prod_{e\in E'}w_e\Bigr)\bigl(\det\mathbf T_{E'}(\mathring\omega)\bigr)^2,
\]
the sum over the sets \(E'\subseteq E_+\) of \(N\) copies, with \(\mathbf T_{E'}\) the square submatrix of the
rows in \(E'\); every term is nonnegative.

Fix a spanning tree, with copy set \(E_{\mathrm{tr}}\subseteq E_+\) of \(N-1\) copies, let
\(e'_1,\ldots,e'_{\mathrm r}\) be the other copies of \(E_+\), with \(\mathrm r=\lvert E_+\rvert-N+1\), and for
each \(k\) let \(p_k\) be the closed walk that traverses \(e'_k\) in its orientation in \(E_+\) and returns
along the path of the tree between its endpoints (\(p_k=(e'_k)\) for a loop copy), with the value
\(\mathring s_k=(\mathbf C\mathring\omega)_{p_k}\) of the cycle map \eqref{eq:link-cycle-map}, the sum of the
generators along \(p_k\) under the reversal convention. The copies \(E'_k=E_{\mathrm{tr}}\cup\{e'_k\}\) form a
connected subgraph with \(N\) vertices and \(N\) copies, the cycle of \(p_k\) with trees attached. While an
attached tree remains, it has a vertex off the cycle that lies on exactly one copy \(e\) of the subgraph;
the column of that vertex has the single nonzero entry \(1\) or \(-e^{\mathring\omega_e}\), in the row of
\(e\), and expanding the determinant along it leaves, up to that entry and a sign, the determinant of the
matrix of the subgraph without that vertex and copy. When only the cycle remains, the row of a copy \(e:j\to i\) that
\(p_k\) traverses against its orientation in \(E_+\) is \(-e^{\mathring\omega_e}\) times the row of the
reversed copy \(\bar e\), which has the entry \(1\) at \(j\) and \(-e^{\mathring\omega_{\bar e}}=-e^{-\mathring\omega_e}\)
at \(i\). After these replacements, and with the \(m\) vertices and the \(m\) copies \(e_a\) of the cycle
ordered along \(p_k\), the matrix has the diagonal entries \(1\) and the entries \(-e^{\mathring\omega_{e_a}}\)
at the positions \((a,a+1)\) modulo \(m\), or the single entry \(1-e^{\mathring\omega_{e'_k}}\) when
\(m=1\). For \(m\geq2\) only the identity and the cyclic shift, of sign \((-1)^{m-1}\), contribute to its
determinant, which is therefore \(1-e^{\mathring s_k}\), as it is for \(m=1\); for \(m\geq2\) its square is
\eqref{eq:glink-cycle-determinant} at \(K=1\) with unit weights. So \(\lvert\det\mathbf T_{E'_k}\rvert\) is
\(\lvert1-e^{\mathring s_k}\rvert\) times a product of factors \(e^{\mathring\omega_e}\) over distinct copies \(e\),
and
\[
e^{-\sum_{e\in E_+}\lvert\mathring\omega_e\rvert}\lvert1-e^{\mathring s_k}\rvert
\leq\lvert\det\mathbf T_{E'_k}(\mathring\omega)\rvert
\leq e^{\sum_{e\in E_+}\lvert\mathring\omega_e\rvert}\lvert1-e^{\mathring s_k}\rvert.
\]
The linear map \(\mathring\omega\mapsto(\mathring s,\mathring\omega_{\mathrm{tr}})\), with
\(\mathring s=(\mathring s_k)_k\in\R^{\mathrm r}\) and \(\mathring\omega_{\mathrm{tr}}=(\mathring\omega_e)_{e\in E_{\mathrm{tr}}}\), is
invertible, because \(p_k\) traverses \(e'_k\) once and no other copy outside the tree, so that
\(\mathring s_k\) is \(\mathring\omega_{e'_k}\) plus a signed sum of tree generators.

(i) Keeping the terms \(E'=E'_k\) of the Cauchy--Binet sum, with \(w_{\min}=\min_ew_e\),
\[
\det\Lambda(\mathring\omega)
\geq w_{\min}^Ne^{-2\sum_e\lvert\mathring\omega_e\rvert}\sum_{k=1}^{\mathrm r}\bigl(1-e^{\mathring s_k}\bigr)^2
\geq w_{\min}^Ne^{-2-2\sum_e\lvert\mathring\omega_e\rvert}\min\bigl(\lvert\mathring s\rvert^2,1\bigr),
\]
because \(\lvert1-e^{\mathrm u}\rvert\geq e^{-1}\min(\lvert\mathrm u\rvert,1)\) for real \(\mathrm u\), by the mean
value theorem for \(\lvert\mathrm u\rvert\leq1\) and by \(\lvert1-e^{\mathrm u}\rvert\geq1-e^{-1}>e^{-1}\)
otherwise, and because \(\sum_k\min(\mathring s_k^2,1)\geq\min(\lvert\mathring s\rvert^2,1)\), the left side
being at least one if some \(\lvert\mathring s_k\rvert\geq1\) and equal to \(\lvert\mathring s\rvert^2\) otherwise.
So \(\Lambda(\mathring\omega)\succ0\) off the subspace \(\{\mathring s=0\}\), which has codimension
\(\mathrm r\geq1\) and is Lebesgue-null, and there
\(\det\Lambda^{-1/2}\leq w_{\min}^{-N/2}e^{1+\sum_e\lvert\mathring\omega_e\rvert}\max(\lvert\mathring s\rvert^{-1},1)\).
Since the map above is a linear isomorphism and \(\nu\) is Gaussian with positive-definite covariance,
there are constants \(\mathrm c_1,\mathrm c_2,\mathrm c_3>0\) such that the density of \(\nu\) is at most
\(\mathrm c_1e^{-\mathrm c_2(\lvert\mathring s\rvert^2+\lvert\mathring\omega_{\mathrm{tr}}\rvert^2)}\) and
\(\sum_e\lvert\mathring\omega_e\rvert\leq\mathrm c_3(\lvert\mathring s\rvert+\lvert\mathring\omega_{\mathrm{tr}}\rvert)\). In
the coordinates \((\mathring s,\mathring\omega_{\mathrm{tr}})\), whose Jacobian is constant,
\(\int\nu(d\mathring\omega)\det\Lambda^{-1/2}\) is therefore at most a constant times
\[
\int_{\R^{N-1}}e^{-\mathrm c_2\lvert\mathring\omega_{\mathrm{tr}}\rvert^2+\mathrm c_3\lvert\mathring\omega_{\mathrm{tr}}\rvert}d\mathring\omega_{\mathrm{tr}}
\int_{\R^{\mathrm r}}e^{-\mathrm c_2\lvert\mathring s\rvert^2+\mathrm c_3\lvert\mathring s\rvert}\max\bigl(\lvert\mathring s\rvert^{-1},1\bigr)d\mathring s,
\]
the first factor being absent when \(N=1\). The first integral is finite, and so is the part of the second
over \(\lvert\mathring s\rvert>1\); on \(\lvert\mathring s\rvert\leq1\) the second integrand is at most
\(e^{\mathrm c_3}\lvert\mathring s\rvert^{-1}\), whose integral there is the area of the unit sphere of
\(\R^{\mathrm r}\) times \(\int_0^1\mathrm u^{\mathrm r-2}d\mathrm u\), finite because \(\mathrm r\geq2\). So
\(\mathfrak Z<\infty\), and \(\mathfrak Z>0\) because its integrand is positive.
(ii) If \(\mathrm r=1\), then \(\lvert E_+\rvert=N\), the Cauchy--Binet sum has the single term
\(E'=E_+=E'_1\), and
\(\det\Lambda\leq\bigl(\prod_ew_e\bigr)e^{2\sum_e\lvert\mathring\omega_e\rvert}(1-e^{\mathring s_1})^2\). On the
bounded set where \(\lvert\mathring s_1\rvert\leq1\) and \(\lvert\mathring\omega_{\mathrm{tr}}\rvert\leq1\), the mean
value theorem gives \(\lvert1-e^{\mathring s_1}\rvert\leq e\lvert\mathring s_1\rvert\), the sum
\(\sum_e\lvert\mathring\omega_e\rvert\) is bounded, and the density of \(\nu\) is bounded below by a positive
constant, so \(\mathfrak Z\) is at least a positive multiple of
\(\int_{-1}^1\lvert\mathrm u\rvert^{-1}d\mathrm u=+\infty\). \hfill$\square$

\propositionheading{Hopping expansion of the induced action}{prop:glink-hopping}
In the setting of \Cref{prop:glink-induced-action} with every \(A_i\succ0\) and \(\mathsf h=0\), let
\(\Lambda^{\mathrm{diag}}(\mathring\omega)=\operatorname{blockdiag}(\Lambda^{\mathrm{diag}}_i)\), where
\(\Lambda^{\mathrm{diag}}_i\) is \(A_i\) plus \(W_e\) for every \(e\in E_+\) ending at \(i\) and
\(e^{\mathring\omega_e^\top}W_ee^{\mathring\omega_e}\) for every \(e\in E_+\) starting at \(i\), a loop counting in
both sums, and let \(\Lambda^{\mathrm{hop}}(\mathring\omega)=\Lambda^{\mathrm{diag}}(\mathring\omega)-\Lambda(\mathring\omega)\).
For a copy \(e\notin E_+\) put \(W_e=e^{\mathring\omega_{\bar e}^\top}W_{\bar e}e^{\mathring\omega_{\bar e}}\), the
reversal law of \Cref{prop:cg-kernel-holonomy}(i).
(i) For every \(\mathring\omega\),
\(\operatorname{spr}((\Lambda^{\mathrm{diag}})^{-1}\Lambda^{\mathrm{hop}})
\leq1-\min_i\min_{z\neq0}z^\top A_iz/z^\top\Lambda^{\mathrm{diag}}_iz<1\), with \(\operatorname{spr}\) the
spectral radius.
(ii) The induced action has the absolutely convergent expansion
\begin{equation}
\mathscr S^{\mathrm{ind}}(\mathring\omega)
=\tfrac12\log\det\Lambda^{\mathrm{diag}}(\mathring\omega)
-\tfrac12\sum_{r\geq1}\frac1r\Tr\bigl[\bigl((\Lambda^{\mathrm{diag}})^{-1}\Lambda^{\mathrm{hop}}\bigr)^r\bigr]-\tfrac{NK}2\log2\pi .
\label{eq:glink-hopping}
\end{equation}
(iii) The \(r\)-th trace is a sum over the closed walks \((e_1,\ldots,e_r)\) of length \(r\), with
\(e_b:i_b\to i_{b-1}\), \(i_r=i_0\), and copies of either orientation, loops included,
\begin{equation}
\Tr\bigl[\bigl((\Lambda^{\mathrm{diag}})^{-1}\Lambda^{\mathrm{hop}}\bigr)^r\bigr]
=\sum_{(e_1,\ldots,e_r)}\Tr\Bigl((\Lambda^{\mathrm{diag}}_{i_0})^{-1}W_{e_1}e^{\mathring\omega_{e_1}}(\Lambda^{\mathrm{diag}}_{i_1})^{-1}W_{e_2}e^{\mathring\omega_{e_2}}\cdots
(\Lambda^{\mathrm{diag}}_{i_{r-1}})^{-1}W_{e_r}e^{\mathring\omega_{e_r}}\Bigr),
\label{eq:glink-walk-sum}
\end{equation}
a weighted Wilson loop: the trace of the reference holonomy
\(e^{\mathring\omega_{e_1}}\cdots e^{\mathring\omega_{e_r}}\) of the walk with a weight inserted at every vertex.
(iv) If \(A_i=a_i\mathsf G\) and \(W_e=w_e\mathsf G\) with scalars \(a_i,w_e>0\), and every link is an
isometry of \(\mathsf G\), then \(\Lambda^{\mathrm{diag}}_i=a^{\mathrm{diag}}_i\mathsf G\), with \(a^{\mathrm{diag}}_i\) equal to
\(a_i\) plus the sum of \(w_e\) over the copies \(e\in E_+\) incident to \(i\), a loop counted twice;
\(\Lambda^{\mathrm{diag}}\) does not depend on the links, and the walk \((e_1,\ldots,e_r)\) contributes
\(\bigl(\prod_bw_{e_b}/a^{\mathrm{diag}}_{i_{b-1}}\bigr)\Tr(e^{\mathring\omega_{e_1}}\cdots e^{\mathring\omega_{e_r}})\). A walk that
reduces to the trivial walk by cancelling backtracks \((e,\bar e)\) contributes a link-independent
constant.
(v) For the triangle, three agents \(i_0,i_1,i_2\) joined by three copies that form one cycle \(p\)
and by no other copy, in the isotropic isometric setting of (iv) and with
\(w^{\mathrm{hop}}=\max_ew_e/\min_ia^{\mathrm{diag}}_i<\tfrac12\),
\begin{equation}
\mathscr S^{\mathrm{ind}}(\mathring\omega)-\mathscr S^{\mathrm{ind}}(0)
=\frac{w_{e_0}w_{e_1}w_{e_2}}{a^{\mathrm{diag}}_{i_0}a^{\mathrm{diag}}_{i_1}a^{\mathrm{diag}}_{i_2}}\mathfrak Q_p(\mathring\omega)+\mathfrak o(\mathring\omega),
\qquad
\lvert\mathfrak o(\mathring\omega)\rvert\leq\frac K5\Bigl(\frac{(2w^{\mathrm{hop}})^5}{1-2w^{\mathrm{hop}}}+\frac{2(w^{\mathrm{hop}})^5}{1-w^{\mathrm{hop}}}\Bigr),
\label{eq:glink-wilson-triangle}
\end{equation}
where \(\mathfrak Q_p=K-\Tr\mathring\Theta_p\) for isometric holonomy.
\status{ESTABLISHED}

\noindent\emph{Proof.}
(i) Each energy term \((z_i-e^{\mathring\omega_e}z_j)^\top W_e(z_i-e^{\mathring\omega_e}z_j)\) contributes
\(W_e\) and \(e^{\mathring\omega_e^\top}W_ee^{\mathring\omega_e}\) to the diagonal blocks, collected in
\(\Lambda^{\mathrm{diag}}\), and \(-W_ee^{\mathring\omega_e}\) and \(-e^{\mathring\omega_e^\top}W_e\) to the blocks
\((i,j)\) and \((j,i)\), so \(\Lambda^{\mathrm{hop}}\) has blocks \(W_ee^{\mathring\omega_e}\) and
\(e^{\mathring\omega_e^\top}W_e\). The same terms with the sign of every link reversed,
\((z_i+e^{\mathring\omega_e}z_j)^\top W_e(z_i+e^{\mathring\omega_e}z_j)\), assemble
\(\Lambda^{\mathrm{diag}}+\Lambda^{\mathrm{hop}}\). Both forms are sums of positive semidefinite terms plus
\(z^\top A_\oplus z\), so \(\Lambda^{\mathrm{diag}}\pm\Lambda^{\mathrm{hop}}\succeq A_\oplus\). With the symmetric
matrix \(\breve\Lambda^{\mathrm{hop}}=(\Lambda^{\mathrm{diag}})^{-1/2}\Lambda^{\mathrm{hop}}(\Lambda^{\mathrm{diag}})^{-1/2}\)
and \(\breve A_\oplus=(\Lambda^{\mathrm{diag}})^{-1/2}A_\oplus(\Lambda^{\mathrm{diag}})^{-1/2}\), this reads
\(-(I-\breve A_\oplus)\preceq\breve\Lambda^{\mathrm{hop}}\preceq I-\breve A_\oplus\), and
\((\Lambda^{\mathrm{diag}})^{-1}\Lambda^{\mathrm{hop}}=(\Lambda^{\mathrm{diag}})^{-1/2}\breve\Lambda^{\mathrm{hop}}(\Lambda^{\mathrm{diag}})^{1/2}\)
has the spectrum of \(\breve\Lambda^{\mathrm{hop}}\). The block-diagonal matrix \(\breve A_\oplus\) has smallest
eigenvalue \(\min_i\min_{z\neq0}z^\top A_iz/z^\top\Lambda^{\mathrm{diag}}_iz\), by the Rayleigh quotient in the
variable \((\Lambda^{\mathrm{diag}}_i)^{1/2}z\), and it lies in \((0,1]\) because
\(\Lambda^{\mathrm{diag}}_i\succeq A_i\succ0\).
(ii) \(\Lambda=(\Lambda^{\mathrm{diag}})^{1/2}(I-\breve\Lambda^{\mathrm{hop}})(\Lambda^{\mathrm{diag}})^{1/2}\), and for
the symmetric matrix \(\breve\Lambda^{\mathrm{hop}}\) with eigenvalues \(\mathring\lambda_k\in(-1,1)\),
\(\log\det(I-\breve\Lambda^{\mathrm{hop}})=\sum_k\log(1-\mathring\lambda_k)=-\sum_{r\geq1}\frac1r\Tr(\breve\Lambda^{\mathrm{hop}})^r\),
absolutely convergent because \(\sum_r\sum_k\lvert\mathring\lambda_k\rvert^r/r<\infty\); similarity gives
\(\Tr(\breve\Lambda^{\mathrm{hop}})^r=\Tr((\Lambda^{\mathrm{diag}})^{-1}\Lambda^{\mathrm{hop}})^r\), and
\eqref{eq:glink-induced-action} at \(\mathsf h=0\) gives \eqref{eq:glink-hopping}.
(iii) The block \((i,j)\) of \((\Lambda^{\mathrm{diag}})^{-1}\Lambda^{\mathrm{hop}}\) is
\((\Lambda^{\mathrm{diag}}_i)^{-1}\) times the sum of \(W_ee^{\mathring\omega_e}\) over the copies \(e:j\to i\) of
either orientation: for \(e\in E_+\) it was computed in (i), and for the reversal of \(e\in E_+\) the block
\(e^{\mathring\omega_e^\top}W_e\) equals \(W_{\bar e}e^{\mathring\omega_{\bar e}}\) with
\(W_{\bar e}=e^{\mathring\omega_e^\top}W_ee^{\mathring\omega_e}\) and \(\mathring\omega_{\bar e}=-\mathring\omega_e\). The
diagonal blocks of the \(r\)-th power are sums over vertex sequences of products of these blocks, which
is \eqref{eq:glink-walk-sum}.
(iv) An isometric link satisfies \(e^{\mathring\omega_e^\top}W_ee^{\mathring\omega_e}=w_e\mathsf G\), and so does
every reversed weight, which gives \(\Lambda^{\mathrm{diag}}_i\) and the factors
\((\Lambda^{\mathrm{diag}}_{i_{b-1}})^{-1}W_{e_b}e^{\mathring\omega_{e_b}}=(w_{e_b}/a^{\mathrm{diag}}_{i_{b-1}})e^{\mathring\omega_{e_b}}\).
A backtrack contributes \(e^{\mathring\omega_e}e^{-\mathring\omega_e}=I\) to the product for every link value.
(v) On the triangle every closed walk winds \(\mathbf k\in\mathbb Z\) times around \(p\), and after
cancelling backtracks its holonomy is conjugate to \(\mathring\Theta_p^{\mathbf k}\), so it contributes its
weight times \(\Tr\mathring\Theta_p^{\mathbf k}\). A walk of length \(r\) has \(r\equiv\mathbf k\pmod2\) and
\(r\geq3\lvert\mathbf k\rvert\), so in \(\mathscr S^{\mathrm{ind}}(\mathring\omega)-\mathscr S^{\mathrm{ind}}(0)\),
where \(\Lambda^{\mathrm{diag}}\) cancels, the walks with \(\mathbf k=0\) cancel, no walk of length four
contributes, and the walks of length three are the six traversals of \(p\), three starting points in two
orientations, all of weight \(w_{e_0}w_{e_1}w_{e_2}/(a^{\mathrm{diag}}_{i_0}a^{\mathrm{diag}}_{i_1}a^{\mathrm{diag}}_{i_2})\). Their term in
\eqref{eq:glink-hopping} is \(-\tfrac12\cdot\tfrac13\cdot6\) times that weight times
\(\Tr\mathring\Theta_p-K\), using
\(\Tr\mathring\Theta_p^{-1}=\Tr(\mathsf G^{-1}\mathring\Theta_p^\top\mathsf G)=\Tr\mathring\Theta_p\) for an
isometry, and for an isometry
\(\tfrac12\langle\mathring\Theta_p-I,\mathring\Theta_p-I\rangle_{\mathsf G}=K-\Tr\mathring\Theta_p\). The
remaining walks have length \(r\geq5\); the triangle has \(2^r+2(-1)^r\leq2^r+2\) closed walks of length
\(r\), the eigenvalues of its adjacency matrix being \(2,-1,-1\), each of weight at most
\((w^{\mathrm{hop}})^r\), and
\(\lvert\Tr\mathring\Theta_p^{\mathbf k}-K\rvert\leq2K\) because an isometry has eigenvalues of modulus one.
Summing \(\tfrac1{2r}(2^r+2)(w^{\mathrm{hop}})^r\cdot2K\leq\tfrac K5((2w^{\mathrm{hop}})^r+2(w^{\mathrm{hop}})^r)\)
over \(r\geq5\) gives the bound. \hfill$\square$

The expansion is the Gaussian-matter form of the induced gauge action of
\citet[Sec.~2, Eqs.~(2)--(4)]{KazakovMigdal1993}, in which integrating out a field coupled to the links
gives one half the logarithm of the determinant of its quadratic operator and an expansion over closed
lattice loops; here the field is the belief fiber, and the leading curvature term is the triangle's
loop defect with the product of the hopping ratios \(w_{e_a}/a^{\mathrm{diag}}_{i_a}\) along the cycle as
coefficient. For the connection Laplacian of a line bundle the determinant, and for a rank-two bundle
with links in \(\mathrm{SL}_2\) its quaternion determinant, has a finite expansion over cycle-rooted
spanning forests \citep[Theorem~1 and Theorem~9]{Kenyon2011VectorBundleLaplacian}; Kenyon lists a combinatorial
interpretation in higher rank as an open question \citep[Sec.~14]{Kenyon2011VectorBundleLaplacian},
and no finite expansion of \(\det\Lambda(\mathring\omega)\) for \(K\geq3\) is stated here.

\section{An explicit spectrum of the linearized link operator}
\label{sec:glink-spectrum}

On a ring of agents with isotropic data, the linearized operator of
\Cref{prop:link-linearized-equation} has an explicit spectrum.

\propositionheading{The linearized operator on a ring of agents and episodes}{prop:glink-linear-spectrum}
Take a ring of \(N\geq2\) agents \(i\in\mathbb Z_N\), with copies \(e_i:i+1\to i\),
\(E_+=\{e_0,\ldots,e_{N-1}\}\), and the closed walk \(p=(e_0,\ldots,e_{N-1})\), which spans the cycle
space of the ring, of dimension one; the learning episodes of \Cref{def:full-graph-episode-joint},
indexed here by \(\mathring n=0,\ldots,\mathring n_1\) with \(\mathring n_1\geq0\); and generator histories
\((\mathring\omega^{(\mathring n)}_{e_i})\) in the history space
\(\R^N\otimes\R^{\mathring n_1+1}\otimes\mathfrak g\) of edge index, episode index, and generator, with the
inner product \(\sum_{\mathring n,i}\langle\cdot,\cdot\rangle_{\mathsf G}\). Let \(\mathrm a>0\) and
\(\mathrm b,\mathrm c,\mathbf m\geq0\), and take the isotropic data of the anchor \(\mathring\alpha=0\) with prior
covariance \(\mathring\Sigma_{e_i}=\mathrm a^{-1}\mathrm{Id}\), the loop coefficient \(\mathrm c\) of \(p\), and a
matter mass \(\mathrm b\) times the identity at every episode, and the increment precision \(\mathbf m\)
between successive episodes, with the history action
\begin{equation}
\mathscr S^{\mathrm{hist}}(\mathring\omega)
=\sum_{\mathring n=0}^{\mathring n_1}\Bigl[\tfrac{\mathrm a+\mathrm b}2\sum_i\bigl\lVert\mathring\omega^{(\mathring n)}_{e_i}\bigr\rVert_{\mathsf G}^2
+\mathrm c\mathfrak Q_p\bigl(\mathring\omega^{(\mathring n)}\bigr)\Bigr]
+\tfrac{\mathbf m}2\sum_{\mathring n=1}^{\mathring n_1}\sum_i\bigl\lVert\mathring\omega^{(\mathring n)}_{e_i}-\mathring\omega^{(\mathring n-1)}_{e_i}\bigr\rVert_{\mathsf G}^2 ,
\label{eq:glink-ring-energy}
\end{equation}
which for \(\mathbf m>0\) is, up to a constant, the negative logarithm of the density of the history
prior of \Cref{def:link-history-prior} with the Gaussian increment as kinetic term, plus the isotropic
second-order matter term. Let \(\mathbf K\) be its Hessian at \(\mathring\omega=0\), the operator of
\Cref{prop:link-linearized-equation} for these data with the matter operator \(\mathrm b\mathrm{Id}\).
(i) As \(\mathring\omega\to0\),
\(\mathfrak Q_p(\mathring\omega)=\tfrac12\lVert\mathbf C\mathring\omega\rVert_{\mathsf G}^2+O(\lVert\mathring\omega\rVert^3)\)
with the cycle map \(\mathbf C\mathring\omega=\sum_i\mathring\omega_{e_i}\), so the loop part of \(\mathbf K\)
is \(\mathrm c\mathbf C^*\mathbf C\), which acts on the edge index as the rank-one matrix
\(\mathrm c\mathbf 1\mathbf 1^\top\) and not as a Laplacian of the ring. (ii) The operator is
\begin{equation}
\mathbf K=\Bigl[(\mathrm a+\mathrm b)I_N\otimes I_{\mathring n_1+1}+\mathrm c\mathbf 1\mathbf 1^\top\otimes I_{\mathring n_1+1}
+\mathbf mI_N\otimes L_{\mathring n_1}\Bigr]\otimes\mathrm{Id}_{\mathfrak g},
\label{eq:glink-ring-operator}
\end{equation}
with \(L_{\mathring n_1}\) the Laplacian of the path on \(\{0,\ldots,\mathring n_1\}\), equal to
\(2\mathrm v_{\mathring n}-\mathrm v_{\mathring n-1}-\mathrm v_{\mathring n+1}\) in the interior and to
\(\mathrm v_0-\mathrm v_1\) and \(\mathrm v_{\mathring n_1}-\mathrm v_{\mathring n_1-1}\) at the ends. Its eigenvalues are
\begin{equation}
\mathrm a+\mathrm b+\mathrm cN+4\mathbf m\sin^2\frac{\pi l}{2(\mathring n_1+1)}
\quad\text{and}\quad
\mathrm a+\mathrm b+4\mathbf m\sin^2\frac{\pi l}{2(\mathring n_1+1)},
\qquad l=0,\ldots,\mathring n_1,
\label{eq:glink-ring-spectrum}
\end{equation}
and for each \(l\), with the episode profile \(\mathrm v^{(l)}\in\R^{\mathring n_1+1}\) of entries
\(\mathrm v^{(l)}_{\mathring n}=\cos(\pi l(\mathring n+\tfrac12)/(\mathring n_1+1))\), \(\mathbf K\) acts by the first
value on the edge-uniform subspace \(\mathbf 1\otimes\mathrm v^{(l)}\otimes\mathfrak g\), of dimension \(K^2\), and by
the second on the zero-edge-sum subspace \(\mathbf 1^\perp\otimes\mathrm v^{(l)}\otimes\mathfrak g\), of dimension
\((N-1)K^2\). These \(2(\mathring n_1+1)\) subspaces are orthogonal and span the history space, so the
multiplicity of an eigenvalue of \(\mathbf K\) is the sum of the dimensions of the subspaces on which
\(\mathbf K\) acts by it, larger than \(K^2\) or \((N-1)K^2\) when values coincide, as when \(\mathrm c=0\) or
\(\mathbf m=0\). The discrete Fourier basis of \(\mathbb Z_N\) is one orthonormal eigenbasis in the edge index.
(iii) Every eigenvalue is at least
\(\mathrm a+\mathrm b\geq\mathrm a>0\), attained by profiles constant in the episode index whose edge sum
vanishes; there is no zero or negative mode. (iv) For \(\mathbf m>0\) and either spatial sector, with
\(\mathrm q\) its constant part \(\mathrm a+\mathrm b+\mathrm cN\) or \(\mathrm a+\mathrm b\), the homogeneous episode
recursion \(\mathbf m(2\mathrm v_{\mathring n}-\mathrm v_{\mathring n-1}-\mathrm v_{\mathring n+1})+\mathrm q\mathrm v_{\mathring n}=0\)
has the characteristic roots
\(1+\tfrac{\mathrm q}{2\mathbf m}\pm\bigl((1+\tfrac{\mathrm q}{2\mathbf m})^2-1\bigr)^{1/2}\), real, positive, and
reciprocal, so its solutions are combinations of one increasing and one decreasing geometric sequence and
none oscillates. \status{ESTABLISHED}

\noindent\emph{Proof.}
In \eqref{eq:link-history-prior} with these data, the base law \(\mathcal N(0,\mathrm a^{-1}\mathrm{Id})\) of each copy
contributes \(\tfrac{\mathrm a}2\lVert\mathring\omega^{(\mathring n)}_{e_i}\rVert_{\mathsf G}^2\) to the negative
logarithm of the density up to a constant, the action contributes \(\mathrm c\mathfrak Q_p\) at every episode,
and the Gaussian increment with mass \(\mathbf m\) contributes the last sum of \eqref{eq:glink-ring-energy},
which gives the identification for \(\mathbf m>0\).
(i) \(e^{\mathbf B}=I+\mathbf B+O(\lVert\mathbf B\rVert^2)\) gives
\(\mathring\Theta_p-I=\sum_i\mathring\omega_{e_i}+O(\lVert\mathring\omega\rVert^2)\), and \(\mathfrak Q_p\) is
half its squared \(\mathsf G\)-norm; the Hessian of
\(\tfrac{\mathrm c}2\lVert\sum_i\mathring\omega_{e_i}\rVert_{\mathsf G}^2\) is \(\mathrm c\mathbf C^*\mathbf C\), which
sends \((\mathring\omega_{e_i})_i\) to the constant profile \(\sum_i\mathring\omega_{e_i}\). (ii) The anchor
and matter terms contribute \((\mathrm a+\mathrm b)\mathrm{Id}\), the loop term
\(\mathrm c\mathbf 1\mathbf 1^\top\) in the edge index at every episode, and the increment term
\(\mathbf mL_{\mathring n_1}\) in the episode index for every edge, which is
\eqref{eq:glink-ring-operator}; these are the summands \(\mathring\Sigma^{-1}\), \(\mathrm c\mathbf C_p^\top\mathbf C_p\),
and \(\mathbf m(-\Delta_n)\) of the operator of \Cref{prop:link-linearized-equation}, with \(\mathrm b\mathrm{Id}\) in
place of its matter operator. The three Kronecker summands commute. The matrix \(\mathbf 1\mathbf 1^\top\)
has eigenvalue \(N\) on \(\mathbf 1\) and \(0\) on \(\mathbf 1^\perp\), and it is circulant, so the Fourier
basis diagonalizes it. For \(\phi=\pi l/(\mathring n_1+1)\) and
\(\mathrm v_{\mathring n}=\mathrm v^{(l)}_{\mathring n}=\cos(\phi(\mathring n+\tfrac12))\), the identity
\(2\cos(\phi(\mathring n+\tfrac12))-\cos(\phi(\mathring n-\tfrac12))-\cos(\phi(\mathring n+\tfrac32))
=4\sin^2(\phi/2)\cos(\phi(\mathring n+\tfrac12))\) gives the interior rows; at \(\mathring n=0\),
\(\cos\tfrac\phi2-\cos\tfrac{3\phi}2=4\sin^2\tfrac\phi2\cos\tfrac\phi2\), and at
\(\mathring n=\mathring n_1\) the same identity holds after the substitution
\(\mathring n\mapsto\mathring n_1-\mathring n\), because \(\phi(\mathring n_1+1)=\pi l\) multiplies both
terms by \((-1)^l\). These \(\mathring n_1+1\) nonzero vectors, with \(\mathrm v^{(l)}_0>0\), are eigenvectors
for distinct eigenvalues of the symmetric \(L_{\mathring n_1}\), hence orthogonal, so they exhaust its
spectrum, and eigenvalues of commuting summands add on the subspaces
\(\mathbf 1\otimes\mathrm v^{(l)}\otimes\mathfrak g\) and \(\mathbf 1^\perp\otimes\mathrm v^{(l)}\otimes\mathfrak g\), which are
orthogonal because \(\mathbf 1\perp\mathbf 1^\perp\) and the \(\mathrm v^{(l)}\) are orthogonal, and whose
dimensions add up to \(N(\mathring n_1+1)K^2\). (iii) Every term of
\eqref{eq:glink-ring-spectrum} beyond \(\mathrm a+\mathrm b\) is nonnegative and vanishes for
\(l=0\) on \(\mathbf 1^\perp\), which is nonzero for \(N\geq2\). (iv) The roots solve
\(\mathrm u^2-(2+\mathrm q/\mathbf m)\mathrm u+1=0\), whose discriminant \((2+\mathrm q/\mathbf m)^2-4\) is positive for
\(\mathrm q>0\); the product of the roots is \(1\) and their sum is positive. \hfill$\square$

The ring is a validation instance, not a model: it shows that the operator of
\Cref{prop:link-linearized-equation} is positive with an explicit gap, that the loop term of a ring is a
single collective mode, and that the episode index enters with the same sign as the anchor, so the most
probable history relaxes monotonically between its ends.

\section{Open problems of the Gaussian link sector}
\label{sec:glink-open}

\openproblemheading{Joint normalization of unanchored matter with latent links}{open:glink-unanchored-normalization}
With positive-definite self terms the joint law of generators and undirected Gaussian matter exists for
every finite base measure (\Cref{prop:glink-induced-action}), and without self terms an unanchored cycle
admits no joint law (\Cref{cor:glink-unanchored-cycle}), because the partition function diverges like
\(\lvert\det(I-\mathring\Theta_p)\rvert^{-1}\). For \(K=1\), with positive weights, zero self terms, and
zero information vector, the partition function of a connected multigraph of cycle rank at least two is
finite under every Gaussian base law with positive-definite covariance
(\Cref{prop:glink-unanchored-scalar}). On a connected multigraph of cycle rank at least two, the
matter precision without self terms is singular exactly when the holonomy group has a nonzero common
fixed vector (\Cref{prop:cg-kernel-holonomy}), and for \(K\geq2\) whether its partition function is
integrable is not settled. It is settled, at a stated \(K\geq2\), for connected multigraphs of cycle rank
at least two with positive-definite weights, zero self terms, zero information vector, and a product
Gaussian base law of the generators with positive-definite covariances, by a proof that
\(\int\nu(d\mathring\omega)\det\Lambda(\mathring\omega)^{-1/2}<\infty\) for every such multigraph and law,
or by a multigraph and law in this class for which the integral diverges.

\openproblemheading{Attainment of the chart infimum without a link prior}{open:glink-chart-infimum}
Without a link prior and for \(K\geq2\), the expected record energy \(\mathbf l\) of
\Cref{prop:glink-chart-infimum-semisimple} has its infimum over the exponential chart at the value of the
rank-one minimizer \(\mathbf M^\star\), unattained, when \(\Tr\mathbf M^\star\geq0\), and strictly above it
when \(\Tr\mathbf M^\star<0\) (\Cref{prop:glink-sample-record}(iii)), where stationary points in the chart can
sit at critical points of the exponential map without being minimizers. In the second case
\Cref{prop:glink-chart-infimum-semisimple} excludes attainment at every chart link whose negative
eigenvalues are semisimple, which covers every chart link for \(K=2\) and \(K=3\) and, for \(K=2\), agrees
with \Cref{cor:link-prior-necessary}(a)(iv). A chart link outside its reach has a negative eigenvalue with a
Jordan block of size at least two, which occurs an even number of times by
\citet[Theorem~1]{culver1966existence}, so that \(K\geq4\); its proof perturbs semisimple negative
eigenvalues to first order and does not reach such a link. It is settled, for \(K\geq4\),
\(W_e,\Sigma_j\succ0\), and data with \(\Tr\mathbf M^\star<0\), by a proof that no chart link having a
negative eigenvalue with a Jordan block of size at least two minimizes \(\mathbf l\) over the chart links,
which with \Cref{prop:glink-chart-infimum-semisimple} gives non-attainment for every \(K\), or by data in
this class and such a chart link that minimizes \(\mathbf l\) over the chart links.
